% TeX root = ../Main.tex

% First argument to \section is the title that will go in the table of contents. Second argument is the title that will be printed on the page.
\section{Mobilní technologie}
\begin{quote}
\textbf{historický vývoj, klasifikace mobilních zařízení – typy, OS, kategorie aplikací, specifické problémy mobilních platforem, senzory a jejich využití v kontextu mobilního uživatele} 
\end{quote}

\subsection{Historický vývoj}

Vývoj mobilních technologií lze rozdělit do několika generací (tzv. ``G'') a klíčových milníků:

\begin{description}
    \item[1G (80. léta)] Analogové sítě (NMT). Sloužily pouze pro přenos hlasu, nízká kvalita, žádná bezpečnost.
    
    \item[2G (90. léta)] Přechod na digitální sítě (GSM). Zavedení SMS, MMS, lepší kvalita hovorů.
    
    \item[3G (počátek 2000+)] Důraz na mobilní internet a multimédia (jednotky Mbps). Video hovory.
\end{description}

\paragraph{Zlomový rok 2007/2008:}

\begin{description}
    \item[2007] Apple představuje iPhone (iOS)~-- revoluce v UI (multitouch, ovládání prsty místo stylusu, absence fyzické klávesnice).
    
    \item[2008] Představen první telefon s OS Android (HTC Dream) a spuštěny obchody s aplikacemi (App Store, Android Market/Google Play).
\end{description}

\begin{description}
    \item[4G/LTE (2010+)] Vysokorychlostní internet. Masivní rozmach streamování videa, cloudových služeb a komplexních aplikací, VoIP.
    
    \item[5G (současnost)] Extrémně nízká latence, vysoká propustnost. Rozvoj IoT (Internet věcí), autonomních vozidel a edge computingu. Nástup skládacích telefonů a integrace AI.
\end{description}

\subsection{Klasifikace mobilních zařízení - typy}

\begin{description}
    \item[Personal Digital Assistants (PDA)] Předchůdci smartphonů. Dotykové obrazovky, virtuální klávesnice, rozpoznávání písma. Organizační funkce: diář, adresář, poznámky. Synchronizace s PC. OS: Palm OS, Windows Mobile.
    
    \item[Smartphony] Nejrozšířenější kategorie, kombinují telefon, počítač a multimediální zařízení. Vlastní operační systém.
    
    \item[Tablety] Větší obrazovka, stejné OS jako smartphony.
    
    \item[Wearables] Chytré hodinky, fitness náramky. Zaměřené na zdraví, notifikace a rychlé interakce.
    
    \item[IoT zařízení] Chytré domácnosti, průmyslové senzory. Zaměřené na automatizaci a sběr dat. Normální operační systémy i RTOS.
\end{description}

\subsection{Operační systémy}

Operační systémy pro mobilní zařízení se liší v přístupu k uživatelskému rozhraní, bezpečnosti, 
vývojářským nástrojům a obchodům s aplikacemi:

\begin{description}
    \item[iOS] Uzavřený ekosystém, vysoká bezpečnost, optimalizace pro Apple hardware. App Store s přísnou kontrolou kvality. Programování v Swiftu a Objective-C.
    
    \item[Android] Otevřený zdrojový kód, linuxové jádro, široká přizpůsobitelnost, velká fragmentace zařízení. Google Play s méně přísnými pravidly. Programování v Javě a Kotlinu.
    
    \item[Další starří OS] Windows Mobile, BlackBerry OS, Symbian, Tizen.
\end{description}

\subsection{Kategorie aplikací}

Aplikace pro mobilní zařízení jsou různého typu, v závislosti na způsobu vývoje a 
cílové platformě:

\begin{description}
    \item[Nativní aplikace] Aplikace napsané pro konkrétní OS (iOS, Android). 
    Vysoký výkon, přístup ke všem funkcím zařízení. Vyžadují vývoj pro každou platformu zvlášť.
    \item[Webové aplikace] Spouštěné v prohlížeči, napsané v HTML, CSS, JavaScriptu. 
    Platformově nezávislé, ale omezený přístup k funkcím zařízení a nižší výkon.
    \item[Hybridní aplikace] Kombinace nativních a webových technologií (např. 
    Apache Cordova, Ionic). Snaží se nabídnout kompromis mezi výkonem 
    a multiplatformní kompatibilitou.
    \item[PWA (Progressive Web Apps)] Webové aplikace s funkcemi blízkými nativním 
    aplikacím (offline režim, push notifikace). Spouštěné v prohlížeči, ale mohou být 
    "instalovány" na zařízení.
    \item[Cross-platformní frameworky] Nástroje pro vývoj jedné kódové základny, 
    která se překládá do nativních aplikací pro různé platformy 
    (např. React Native, Flutter). Snaží se minimalizovat potřebu psaní specifického 
    kódu pro každou platformu.
\end{description}

\subsection{Specifické problémy mobilních platforem}
\begin{description}
    \item[Fragmentace] Různé verze OS, různé hardwarové konfigurace. Ztížená 
    optimalizace a testování aplikací.
    \item[Bezpečnost] Neautorizovaný přístup k datům, malware (zejména u Androidu).
    Ochrana osobních údajů, sandboxing, oprávnění aplikací.
    \item[Omezené zdroje] Nižší výkon, omezená baterie, menší úložiště. Nutnost 
    optimalizace kódu a efektivního využití zdrojů.
    \item[Uživatelské rozhraní] Menší obrazovka, dotykové ovládání. Nutnost 
    přizpůsobení UI pro různé velikosti a rozlišení.
    \item[Síťová konektivita] Nestabilní připojení, přechod mezi sítěmi (Wi-Fi, 4G/5G). 
    Práce s offline režimem. Snížená spotřeba dat.
\end{description}

\subsection{Senzory a jejich využití v kontextu mobilního uživatele}

Oproti desktopovým počítačům mají mobilní zařízení širokou škálu integrovaných senzorů, 
které umožňují nové způsoby interakce s okolním světem:

\begin{description}
    \label{des:senzory}
    \item[Akcelerometr] Měří zrychlení, detekce pohybu a orientace zařízení. 
    Používá se pro krokoměry, detekci pádu, ovládání gesty.
    \item[Gyroskop] Měří rychlost rotace zařízení. Používá se pro herní aplikace, 
    navigaci a realitu.
    \item[GPS] Poskytuje geografické informace. Používá se pro navigaci, mapy a 
    lokální služby.
    \item[Proximity senzor] Detekuje přítomnost předmětů v blízkosti zařízení. 
    Používá se pro automatické vypnutí obrazovky při volání.
    \item[Světelný senzor] Měří úroveň osvětlení v okolí zařízení. Používá se pro 
    automatickou úpravu jasu obrazovky.
    \item[Biometrické senzory] Čtečky otisků prstů, rozpoznávání obličeje. Používají 
    se pro autentizaci a zabezpečení.
    \item[Magnetometr] Měří magnetické pole, funguje jako kompas. Používá se pro 
    navigaci a orientaci.
\end{description}

\hyperref[sub:vsHW]{Zpět} k využití specifického HW (Kapitola 2.)

\section{Vývoj mobilních aplikací}
\begin{quote}
\textbf{specifika vývoje mobilního softwaru, nativní vs hybridní vývoj, Základní charakteristika OS Android (architektura, nástroje) nebo iOS. Vývojářské nástroje a jazyky/frameworky pro nativní vývoj. Obecná problematika tvorby UI (Views, Compose, Swift UI případně dalších). Využití specifického HW, způsoby napojení na IS}
\end{quote}

\subsection{Specifika vývoje mobilního softwaru}

Vývoj mobilních aplikací má svá specifika, která se liší od vývoje pro desktopové nebo 
webové platformy. Klíčové aspekty zahrnují:

\begin{description}
    \item[Optimalizace pro omezené zdroje] Mobilní zařízení mají omezenou výpočetní 
    kapacitu, paměť a baterii. Vývojáři musí optimalizovat kód a minimalizovat 
    spotřebu zdrojů.
    \item[Uživatelské rozhraní] Menší obrazovka a dotykové ovládání vyžadují 
    přizpůsobení UI pro různé velikosti a rozlišení. Důraz na jednoduchost a 
    intuitivnost.
    \item[Fragmentace] Různé verze OS, různé hardwarové konfigurace ztěžují 
    optimalizaci a testování aplikací.
    \item[Bezpečnost] Mobilní platformy jsou náchylné k bezpečnostním hrozbám, 
    jako je malware a neoprávněný přístup k datům. 
    Vývojáři musí implementovat bezpečnostní opatření.
    \item[Integrace s hardwarem] Mobilní zařízení mají širokou škálu senzorů a 
    funkcí (GPS, kamera, akcelerometr).
    \item[Distribuce a monetizace] Mobilní aplikace jsou distribuovány prostřednictvím 
    obchodů s aplikacemi (App Store, Google Play), které
    mají specifické požadavky a pravidla. 
    \item[Sandboxing a oprávnění] Mobilní platformy používají sandboxing pro izolaci 
    aplikací a systém oprávnění pro kontrolu přístupu k funkcím zařízení.
\end{description}

\subsection{Nativní vs. hybridní vývoj}

Nativní vývoj znamená psaní aplikace specificky pro jeden operační systém (iOS, Android), 
zatímco hybridní vývoj využívá webové technologie (HTML, CSS, Javascript) a běží v 
rámci WebView nebo nativního kontejneru.

Nativní vývoj nabízí lepší výkon, přístup ke všem funkcím zařízení a lepší uživatelský 
zážitek, ale vyžaduje vývoj pro každou platformu zvlášť. Hybridní vývoj umožňuje 
multiplatformní kompatibilitu s jednou codebase, ale může mít nižší výkon a omezený 
přístup k funkcím zařízení.

\begin{description}
    \item[Nativní vývoj] iOS (Swift, Objective-C), Android (Java, Kotlin). 
    \item[Hybridní vývoj] Frameworky jako Apache Cordova, Ionic.
\end{description}

\subsection{Základní charakteristika OS Android}

Android je otevřený operační systém založený na linuxovém jádře, spravovaný společností 
Google. Nad jadrem je vrstva Android Runtime (dříve Dalvik) pro spouštění aplikací, 
knihovny pro různé funkce a aplikační rámec pro vývojáře. Aplikace jsou psány v Javě 
nebo Kotlinu a běží v sandboxu (izolovaně) pro bezpečnost. 
Komponenty Androidu:
\begin{description}
    \item[Activity] Uživatelské rozhraní a interakce.
    \item[Service] Běží na pozadí pro dlouhodobé operace.
    \item[Broadcast Receiver] Reaguje na systémové události. (např. změna sítě)
    \item[Content Provider] Sdílení dat mezi aplikacemi.
\end{description}
Vývojové nástroje zahrnují Android Studio, SDK, emulátory. Způsoby tvorby UI: 
\begin{description}
    \item[XML views] Tradiční přístup k tvorbě UI pomocí XML a Java/Kotlin.
    \item[Jetpack Compose] Moderní deklarativní framework pro tvorbu UI v Kotlinu.
\end{description}
Struktura OS:
\begin{description}
    \item[Jádro] Linuxové jádro pro správu hardwaru a systémových zdrojů.
    \item[Hal] Hardware Abstraction Layer pro komunikaci s hardwarem.
    \item[Android Runtime] Spouštění aplikací, správa paměti.
    \item[Knihovny] Nativní knihovny pro grafiku, databáze, síť, atd.
    \item[Java API] Pro vývoj aplikací.
    \item[Aplikační rámec] API pro vývojáře pro přístup k funkcím zařízení.
    \item[Aplikace] Uživatelské aplikace a systémové aplikace.
\end{description}

\subsection{Základní charakteristika OS iOS (Volitelné stačí znát pouze Android)}
iOS je uzavřený operační systém vyvinutý společností Apple pro své mobilní zařízení.
Je založen na unixovém jádře a používá architekturu MVC (Model
-View-Controller). Aplikace jsou psány v Swiftu nebo Objective-C a běží v sandboxu 
pro bezpečnost.
Vývojové nástroje zahrnují Xcode, iOS SDK, simulátory (Apple nazev pro emulátory). Způsoby tvorby UI:
\begin{description}
    \item[Storyboardy] Vizualní nástroj pro návrh UI a navigace mezi obrazovkami.
    \item[SwiftUI] Moderní deklarativní framework pro tvorbu UI v Swiftu.
\end{description}
Struktura OS:
\begin{description}
    \item[Jádro] Unixové jádro pro správu hardwaru a systémových zdrojů.
    \item[Core OS] Nízká úroveň služeb, jako správa paměti, bezpečnost, atd.
    \item[Core Services] Služby pro síť, data, atd.
    \item[Media] Knihovny pro grafiku, zvuk, video.
    \item[Application Framework] API pro vývojáře pro přístup k funkcím zařízení.
    \item[Aplikace] Uživatelské aplikace a systémové aplikace.
\end{description}

\subsection{Vývojářské nástroje a jazyky/frameworky pro nativní vývoj}
\begin{description}
    \item[iOS] Xcode IDE, Swift a Objective-C jazyky, iOS SDK, Interface Builder (součást Xcode) pro návrh UI. Framework SwiftUI pro deklarativní UI.
    \item[Android] Android Studio IDE, Java a Kotlin jazyky, Android SDK, XML pro návrh UI, Framework Jetpack Compose pro
    deklarativní UI.
\end{description}

\subsection{Obecná problematika tvorby UI}
Tvorba uživatelského rozhraní pro mobilní zařízení zahrnuje několik klíčových aspektů:
\begin{description}
    \item[Responsivita] UI musí být přizpůsobeno různým velikostím obrazovek a 
    rozlišením. Použití flexibilních layoutů a komponent.
    \item[Dotykové ovládání] UI musí být optimalizováno pro dotykové ovládání, s 
    dostatečně velkými tlačítky a intuitivními gesty.
    \item[Výkon] UI by mělo být plynulé a rychlé, minimalizovat zpoždění a 
    optimalizovat načítání dat.
    \item[Designové principy] Dodržování platformových designových směrnic 
    (Material Design pro Android, Human Interface Guidelines pro iOS) pro konzistentní 
    a intuitivní uživatelský zážitek.
    \item[Testování] Testování UI na různých zařízeních a verzích OS pro zajištění 
    kompatibility a správné funkčnosti.
\end{description}

\paragraph*{Views}
Tradiční přístup k tvorbě UI pomocí XML (Android) nebo Storyboardů (iOS) a 
programování v Java/Kotlin nebo Swift/Objective-C. Umožňuje detailní kontrolu nad 
vzhledem a chováním UI, ale je náročnější na údržbu a aktualizace.
XML views v Androidu definují rozložení a v kódu se odkazuje pomocí ID. Komponenty
jako: TextView, Button, ImageView, atd.

\paragraph*{Compose a SwiftUI}
Moderní deklarativní frameworky pro tvorbu UI, které umožňují psaní UI kódu přímo v 
programovacím jazyce (Kotlin pro Jetpack Compose a Swift pro SwiftUI). Snadné 
přizpůsobení různým typům zařízení, lepší integrace s logikou aplikace a jednodušší 
aktualizace UI na základě změn stavu.
Jetpack Compose tvoří UI pomocí funkcí s anotací @Composable, které vracejí UI 
komponenty. 
\paragraph*{Architektury:} 
\begin{description}
    \item[MVVM] Model-View-ViewModel, ViewModel obsahuje data pro UI. Model komunukuje s UI a dodává data pro zobrazení. 
    ViewModel zpracovává logiku a aktualizuje UI na základě změn stavu.
    \item[MVI] Model-View-Intent, UI pouze volá akce (intenty), které mění stav modelu, a UI se aktualizuje na základě změn stavu.
\end{description}

\subsection{Využití specifického HW} 
\label{sub:vsHW}
Mobilní zařízení mají širokou škálu integrovaných senzorů a funkcí, které umožňují 
nové způsoby interakce s okolním světem. Vývojáři mohou využít tyto senzory pro různé 
účely, jako je detekce pohybu, navigace, autentizace a další.

\begin{description}
    \item[Senzory] Akcelerometr, gyroskop, GPS, proximity senzor, světelný senzor, 
    biometrické senzory, magnetometr. Viz. \hyperref[des:senzory]{Senzory}
    \item[Kamera] Využití pro fotografování, skenování QR kódů, rozšířenou realitu.
    \item[Mikrofon] Využití pro hlasové ovládání, rozpoznávání řeči, nahrávání zvuku.
    \item[Bluetooth a NFC] Využití pro komunikaci s jinými zařízeními, 
    bezkontaktní platby, připojení periferií.
\end{description}
Přístup k těmto funkcím se obvykle realizuje prostřednictvím API poskytovaných 
operačním systémem, je nutné získat oprávnění od uživatele pro přístup k citlivým 
funkcím a datům. Od androidu 6.0 a iOS 10 je nutné žádat o oprávnění za běhu 
aplikace než pouze při instalaci.

\subsection{Způsoby napojenína IS}
Mobilní aplikace často potřebují komunikovat s informačními systémy (IS) pro získávání 
dat, autentizaci uživatelů, synchronizaci a další funkce. Způsoby napojení na 
IS zahrnují:
\begin{description}
    \item[REST API] Nejčastější způsob komunikace, využívá HTTP protokol pro 
    výměnu dat ve formátu JSON nebo XML.
    \item[GraphQL] Alternativa k REST, umožňuje klientům specifikovat přesně, jaká data 
    potřebují. Protokol: HTTP
    \item[WebSocket] Pro real-time komunikaci, umožňuje obousměrnou komunikaci mezi 
    klientem a serverem.
\end{description}
Při napojení na IS je důležité zohlednit bezpečnostní aspekty, jako je autentizace, 
autorizace, šifrování dat a ochrana proti útokům (např. SQL injection, XSS).

\paragraph*{Implementace na Androidu:} Knihovny jako Retrofit a OkHttp (starší) pro REST API a WebSocket, 
Apollo pro GraphQL. Autentizace může být řešena pomocí OAuth 2.0, JWT, nebo API klíčů. 
Lokální ukládání dat může být realizováno pomocí SQLite (knihovna Room), 
SharedPreferences nebo DataStore. Ke komunikaci Jetpack Compose využívá
Interceory: zachytávají události, Authenticator specialní Inreceptor pro autentizaci, 
Logging Interceptor pro logování síťové komunikace.

\newpage

\section{Inversion of Control (IoC)}
\begin{quote}
\textbf{principy, význam, a použití přístupu, vztah IoC a Dependency Injection. Příklady použití, nástroje podporující IoC}
\end{quote}

\subsection{Principy}
Inversion of Control (IoC) je princip návrhu softwaru, který obrací tradiční řízení 
toku programu. Místo toho, aby komponenty samy řídily své závislosti a životní cyklus, 
IoC umožňuje externímu kontejneru nebo frameworku řídit tyto aspekty.
IoC je všeobecný koncept, který může být implementován různými způsoby, z nichž 
nejběžnější je Dependency Injection (DI).

\subsection{Význam}
\begin{description}
    \item[Modularita] Komponenty jsou méně závislé na konkrétních implementacích 
    svých závislostí, což usnadňuje jejich výměnu a údržbu.
    \item[Testovatelnost] Snadnější testování díky možnosti snadno nahradit závislosti 
    mocky.
    \item[Flexibilita] Snadnější změna implementace závislostí bez nutnosti měnit kód, který je používá.
    \item[Znovupoužitelnost] Komponenty mohou být znovu použity v různých kontextech díky oddělení od konkrétních implementací závislostí.
\end{description}

\subsection{Použití přístupu}
IoC je možné implementovat několika přístupy.
\begin{description}
    \item[Dependency Injection (DI)] Závislosti jsou "injektovány" do komponenty zvenčí, obvykle prostřednictvím konstruktoru, setterů nebo rozhraní.
    \begin{verbatim}
    public ServiceA(ServiceB serviceB) {
        this.serviceB = serviceB;
    }
    \end{verbatim}
    \item[Service Locator] Komponenta získává své závislosti z centrálního registru nebo "služby", která je zodpovědná za poskytování instancí.
    \begin{verbatim}
    ServiceB serviceB = ServiceLocator.getService(ServiceB.class);
    \end{verbatim}
    \item[Event-driven] Komponenty komunikují prostřednictvím událostí, které jsou zpracovávány externím systémem.
    \begin{verbatim}
        eventBus.post(() => {
            // Handle event
        });
    \end{verbatim}
\end{description}


\subsection{Vztah IoC a Dependency Injection}
Dependency Injection (DI) je konkrétní implementace principu Inversion of Control (IoC).
Zatímco IoC je obecný koncept, který může být realizován různými způsoby, 
DI je konkrétní způsob, jak dosáhnout IoC tím, že umožňuje externímu systému 
"injektovat" závislosti do komponenty.
DI je nejběžnější a nejpoužívanější způsob implementace IoC.

\subsection{Příklady použití a nástroje podporující IoC}
IoC se používá v mnoha moderních frameworkách a knihovnách pro vývoj softwaru, 
zejména v oblasti webových a mobilních aplikací. Například:
\begin{description}
    \item[Spring Framework (Java)] Používá IoC kontejner pro správu závislostí a životního cyklu komponent.
    Využití DI pomocí anotací jako @Autowired pro automatické vkládání závislostí. @Component pro označení komponenty, 
    @Service pro označení služby, @Repository pro označení repozitáře.
    Princip: Dependency Injection (DI) pomocí anotací. 
     \item[Angular (TypeScript)] Používá DI pro správu služeb a komponent.
     Princip: Dependency Injection (DI) pomocí dekorátorů jako @Injectable.
     \item[ASP.NET Core (C\#)] Podporuje DI pro správu závislostí v aplikacích.
     Princip: Dependency Injection (DI) pomocí vestavěného kontejneru a metod jako AddTransient, AddScoped, AddSingleton.
     \item[NestJS (Node.js)] Používá DI pro správu závislostí v serverových aplikacích.
     Princip: Dependency Injection (DI) pomocí dekorátorů jako @Injectable a modul
     \item[Android] Knihovny jako Dagger a Hilt pro implementaci DI v Android aplikacích.
        Princip: Dependency Injection (DI) pomocí anotací @Inject a generování kódu pro správu závislostí.
\end{description}

\newpage

\section{Architektura informačních systémů}
\begin{quote}
\textbf{vrstva business logiky, prezentační vrstva, datová vrstva - využití relačních a nerelačních DB v kontextu IS, návrhové vzory pro IS, metody tvorby, komunikace mezi IS, automatizace, DevOps, využití do Cloudu}
\end{quote}
Informační systém je softwarová aplikace zpracovávající a poskytující 
data pro podporu plánování, rozhodování a řízení. Je to celek složený z 
počítačového hardwaru a souvisejícího softwaru, k němuž patří také lidé, 
kteří tento hardware a software využívají, a procesy (činnosti), které 
přitom vykonávají za účelem sběru, zpracování a šíření informací 
potřebných k plánování, rozhodování a řízení.

\paragraph*{Charakteristika IS:}
\begin{itemize}
    \item Aplikace nebo integrovaná sada aplikací
    \item Slouží pro celou organizaci (ne jen pro jednoho uživatele)
    \item Sdílení dat mezi uživateli
    \item Běží mimo počítače uživatelů (na centrální infrastruktuře/v cloudu)
    \item Uživatelé využívají pro přístup k IS svých klientských zařízení
\end{itemize}

\paragraph*{Typy IS:}
\begin{description}
    \item[Podnikové IS] -- Systémy využívané komerčními organizacemi a dostupné zejména pro jejich pracovníky. (Univerzální, Speciální, Na míru)
    \item[Veřejné IS] -- Veřejně či komunitně dostupné IS. (Encyklopedie (např. wikipedie), knihovní systémy, vyhledávače, mapové systémy, …)
\end{description}

\subsection{Vrstvy IS}
\label{sub:vrstvy}
\paragraph*{Vrstva business logiky:}
Vrstva business logiky koordinuje tok dat mezi prezentační a datovou vrstvu. Obsahuje 
většinu algoritmů systému. Většinou reprezentovaná třídami v objektově 
orientovaném jazyce (Java, C\#, JavaScript, …). S okolím komunikuje pomocí 
rozhraní, např. REST API, SOAP, atd.

\paragraph*{Vrstva prezentační:}
Prezenční vrstva zpracovává vstupy od uživatele a volá služby vrstvy 
business logiky. Sama neobsahuje žádnou business logiku, kromě validace, 
jednoduché výpočty (důraz na efektivitu a výkon).
Často postavená na návrhovém vzoru Model-View-Controller a jeho 
variacích. Časté využití technologií založených na JavaScriptu. (Html, CSS, JS, Reacts, Flutter, WPF)

\paragraph*{Vrstva datová:}
Datová vrstva obsehuje služby pro ukládání, načítání dat a dotazování.
Pracuje s databázemi, cache, soubory, atd. Primárně zajiťuje persistentní 
uložení dat, jejich dostupnost a integritu.
Pojmy:
\begin{itemize}
    \item \textbf{CRUD} -- Create, Read, Update, Delete.
    \item \textbf{ORM} -- Mapování objektové struktury aplikace na relační strukturu databáze.
    \item \textbf{Scéma} -- Schéma je datový model, který popisuje data a vztahy mezi nimi. (datové typy, domény - povolené/povinné hodnoty, relace - vazby mezi entitami)
\end{itemize}
Většinou reprezentovaná databází:
\begin{itemize}
    \item \textbf{SQL:} Mají dané schéma. Jsou reprezentované relačním 
    datovým modelem. Důraz na integritu dat. Princip ACID:
    \begin{itemize}
        \item Atomicity -- Pokud se transakce skládá z několika kroků, databáze provede všechny, nebo žádný. Pokud v průběhu dojde k chybě, celá transakce se vrátí zpět (tzv. rollback).
        \item Consistency -- Transakce musí převést databázi z jednoho platného stavu do druhého. Veškerá data musí splňovat definovaná pravidla, omezení (constraints) a integritu.
        \item Isolation -- Souběžně probíhající transakce se vzájemně neovlivňují. Výsledky rozpracované transakce nejsou viditelné pro ostatní procesy, dokud není transakce zcela dokončena.
        \item Durability -- Po úspěšném dokončení transakce (commit) jsou změny trvale zapsány a nevymaže je ani náhlý výpadek proudu či pád systému.
    \end{itemize}
    Vhodné pro IS s komplexními datovými vztahy (ERP, Účetní systémy apod).
	Možné pouze vertikální škálování (zvýšení výkonu jednoho stroje).
    Technologie například MySQL, PostgreSQL, MS SQL, Oracle, MariaDB.
    \item \textbf{NOSQL:} Mají flexibilní schema. Důraz na distribuci dat a škálovatelnost.
    Typy:
    \begin{itemize}
        \item Dokumentové -- Ukládání ve formátu JSON/BSON. Např. MongoDB, CouchDb.
        \item Key-Value -- Ukládání do mapy zvané Hash. Často jako cache. Např. Redis.
        \item Sloupcové -- Čtení po sloupcích což je efektivní pro velké datasety. Např. Cassandra.
        \item Grafové -- Specializace na propojená data. Využití například u sociálních sítí. Např. Neo4j.
    \end{itemize}
    Vhodné pro BigData, aplikace s rychlým vývoje struktura, vysoká dostupnost, realtime (chat, senzory, IoT).
    Podpora transakcí je slabší nebo žádná. Nemají standardizovaný dotazovací
    jazyk jako SQL. Jsou vhodné na horizontální škálování (přidání dalších strojů).
  
    Více k \hyperref[sec:nosql]{NOSQL}. (Kapitola 6.)
    \label{sub:dv}

\end{itemize}

\subsection{Architektury IS}
Architektura je sadou zásadních rozhodnutí týkajících se struktury 
systému, návrhu komponent, jejich vzájemných vztahů a vztahů s okolím.
Jde o nejvyšší (a nejdůležitější) úroveň SW návrhu.

\paragraph*{Co řeší:}
\begin{itemize}
    \item Rozčlenění systému na samostatné celky (subsystémy/komponenty, mikroslužby)
    \item Hlavní komponenty IS a jejich rozhraní
    \item Hlavní vrstvy systému, způsob jejich realizace a vzájemné komunikace
    \item Doménový model
    \item Použité technologie
    \item Způsob integrace s okolními systémy a periferiemi
\end{itemize}
\paragraph*{Typy:}
\begin{description}
    \item[Vrstvená architektura:] Také označovaná jako „Vrstvený monolit“ nebo jen „Monolit“.
    Systém je vytvořen jako jeden nasaditelný celek, který je 
    vnitřně rozčleněn na vrstvy, komponenty. Obsahuje všechny zmíněné \hyperref[sub:vrstvy]{vrstvy}.
    Snadné ladění a testování ale špatná údržba a škálovatelnost.
    \item[Modulární monolit:] Architektura je rozdělená do modulů, což zlepšuje
    organizaci, ale stále běží jako jedna aplikace. 
    (Modul = samostatná logická jednotka, např. fakturace, objednávky)
    \item[Architektura založená na mikroslužbách:] Rozdělení jednoho 
    systému na více samostatných a samostatně nasaditelných celků.
    (Mikroslužba = samostatně nasaditelná komponenta IS, často se samostatnou DB, která se zaměřuje na omezenou oblast business logiky)
    Dobrá škálovatelnost,dostupnost a podpora paralelního vývoje.
    Složité nasazení, údržba a ladění.
    Služby jsou spravovány nástrojem pro orchestraci, například 
    Kubernetes nebo Docker swarm.
    \label{sub:mikrosluzby}
    \hyperref[sub:mikrosluzby2]{Zpět} (Kapitola 5.)
\end{description}


\paragraph*{Konkrétní návrhové vzory:}
\begin{description}
    \item[MVC (Model-View-Controller):] Oddělení datového modelu (Model), uživatelského rozhraní (View) a řídicí logiky (Controller). Cílem je, aby změna v UI minimálně ovlivňovala datovou logiku a naopak.

    \item[MVVM (Model-View-ViewModel):] Používaný ve webových a mobilních aplikacích (Angular, WPF). ViewModel připravuje data z modelu přímo pro potřeby zobrazení a využívá obousměrnou vazbu dat (data binding).

    \item[Hexagonální architektura (Ports and Adapters):] Zaměřuje se na izolaci jádra business logiky od vnějších vlivů (databáze, API, UI). Jádro definuje "porty" (rozhraní) a okolí implementuje "adaptéry". Díky tomu lze snadno vyměnit např. typ databáze bez zásahu do logiky.
\end{description}

\subsection{Metody tvorby IS}
Metodiky se dnes dělí primárně podle přístupu k plánování a změnám na 
tradiční (prediktivní) a agilní. Většinou se kombinují.

\subsubsection{Tradiční (prediktivní) přístupy}
Tyto metody předpokládají, že lze na začátku projektu definovat všechny požadavky, které se pak v průběhu vývoje nemění.

\begin{description}
    \item[Vodopádový model (Waterfall):] Lineární proces, kde každá fáze (Požadavky, Návrh, Implementace, Testování, Nasazení, Údržba) následuje až po úplném dokončení té předchozí. Těžkopádný přístup.
    \item[RUP (Rational Unified Process):] Robustní metodika, která je sice iterativní, ale stále velmi formální a orientovaná na detailní dokumentaci a UML modely.
    Iterativní metoda se čtyřmi fázemi:
    \begin{itemize}
        \item Zahájení (ujasnit vizi a cíle)
        \item Upřesnění (návrh systému)
        \item Konstrukce (vývoj, testování)
        \item Přechod do provozu
    \end{itemize}
    \item[PRINCE2 (PRojects IN Controlled Environment):] Jde o 
    projektový rámec, který dává jasné role, fáze a rozhodovací 
    mechanismy. Dá se použít na velký projekt i malou aplikaci.Projekt 
    dělí na řiditelné etapy s měřitelnými výstupy.
    \begin{figure}[H]
    \centering
    \includegraphics[width=0.5\textwidth]{prince2.png}
    \caption{PRINCE2 - Aspekty řízení projektu}
    \end{figure}
\end{description}

\subsubsection{Agilní přístupy}
Reagují na fakt, že požadavky zákazníka se v čase mění. Namísto velkého plánu na začátku sázejí na krátké cykly a průběžné dodávání funkčního softwaru.

\begin{description}
    \item[SCRUM:] Nejpoužívanější agilní rámec. Práce probíhá ve sprintech (2-4 týdny). Klíčové jsou role (Scrum Master, Product Owner) a pravidelné schůzky (Daily Stand-up, Retrospektiva).
    Velice flexibilní, snadné změny během vývoje a rychlá zpětná vazba.
    \item[KANBAN:] Zaměřuje se na plynulost toku práce (flow) a vizualizaci úkolů na tabuli. Neřeší pevné sprinty, ale soustředí se na omezení rozdělané práce (WIP limit).
    \item[XP (Extrémní programování):] Zaměřeno na kvalitu kódu a technické praktiky jako párové programování, TDD (vývoj řízený testy) a časté vydávání verzí.
\end{description}

\subsection{Komunikace mezi IS}
Způsob, jakým si systémy vyměňují data, zásadně ovlivňuje jejich propustnost, stabilitu a míru provázanosti (coupling).

\paragraph*{1. Synchronní komunikace (Request-Response)}
Klient odešle požadavek a čeká na odpověď serveru. Pokud server neodpovídá, klient je blokován.

\begin{description}
    \item[REST (Representational State Transfer):] Dnes nejrozšířenější standard využívající protokol HTTP a formát JSON. Je bezstavový a využívá standardní metody (GET, POST, PUT, DELETE).
    \item[SOAP (Simple Object Access Protocol):] Starší, robustnější protokol využívající XML. Často se používá v bankovnictví nebo enterprise sféře díky silnému typování a bezpečnosti (WS-Security).
    \item[gRPC:] Moderní, vysoce výkonný framework od Google. Používá binární formát Protocol Buffers a HTTP/2. Ideální pro interní komunikaci mezi mikroslužbami díky nízké latenci.
\end{description}

\paragraph*{2. Asynchronní komunikace (Messaging)}
Systémy spolu komunikují nepřímo skrze prostředníka (Message Broker). Odesílatel nečeká na zpracování, zpráva je uložena do fronty.

\begin{description}
    \item[Message Queues (Bod-bod):] Zpráva je doručena právě jednomu příjemci. Pokud je příjemců více, rozdělují si práci mezi sebou (např. RabbitMQ).
    \item[Publish/Subscribe (Pub/Sub):] Odesílatel publikuje zprávu do tématu (Topic) a tu obdrží všichni odběratelé, které dané téma zajímá. Typickým zástupcem je Apache Kafka.
    \item[Výhody:] Časové rozdělení (systémy nemusí běžet současně), odolnost proti špičkám v zátěži, snadné škálování.
\end{description}

\paragraph*{3. Ostatní formy integrace}
\begin{description}
    \item[Webhooks:] "Obrácené API". Server pošle HTTP požadavek na předem definovanou URL klienta, jakmile nastane určitá událost (např. proběhnutí platby).
    \item[Sdílená databáze:] Více systémů přistupuje ke stejné DB. Považováno za antipattern, protože změna schématu v jednom systému rozbije ostatní.
    \item[Souborový přenos (Batch processing):] Předávání dat v dávkách (např. CSV nebo XML soubory přes SFTP). Často u starších systémů (legacy) nebo pro přenos velkých objemů dat.
    \item[WebSockets:] Trvalé spojené mezi klientem a serverem. Vhodné pro realtime app (chat, hry).
    \item[Enterprice Service Bus (ESB):] Middleware plaforma pro 
    centralizovanou integraci. Podporuje různé protokoly. Např. Apache Came, WSO2.
\end{description}

\subsection{Automatizace a DevOps}
DevOps (spojení slov Development a Operations) je metodika, která bourá bariéry mezi vývojáři a IT správci.
Automatizace je zahrnuje build, testování, nasazení, monitoring atd.

\paragraph*{CI/CD (Continuous Integration / Continuous Deployment):}
Odstraňuje manuální kroky při nasazování a zajišťuje tak automatizaci.
\begin{itemize}
    \item \textbf{CI (Průběžná integrace):} Automatizovaný proces, kdy se kód od vývojářů pravidelně slučuje do hlavní větve. Každý "commit" spustí automatické sestavení (build), sadu testů, analýza kódu, a další úkony. Proces, který se snaží udržet všechny (vývojáře, testery …) synchronizovány,
    \item \textbf{CD (Průběžné doručování/nasazování):} Navazuje na CI. Automaticky připraví aplikaci k nasazení. U \textit{Deploymentu} jde kód po úspěšných testech rovnouna prostředí bez lidského zásahu.
\end{itemize}

\paragraph*{Klíčové praktiky a nástroje:}
\begin{itemize}
    \item \textbf{IaC (Infrastructure as Code):} Správa serverů a sítí pomocí konfiguračních souborů místo manuálního nastavování (např. Terraform, Ansible). Infrastruktura se pak chová jako software (lze ji verzovat).
    \item \textbf{Kontejnerizace:} Balení aplikace a všech jejích závislostí do jednoho balíčku (kontejneru), který běží stejně na počítači vývojáře i na serveru. Dominantním nástrojem je \textbf{Docker}.
    \item \textbf{Orchestrace:} Správa velkého množství kontejnerů, jejich škálování a restartování v případě výpadku. Standardem je \textbf{Kubernetes}.
    \item \textbf{Monitoring a Logging:} Neustálé sledování výkonu systému a sběr chybových logů (např. ELK stack - Elasticsearch, Logstash, Kibana), aby bylo možné reagovat na problémy dříve, než si jich všimne uživatel.
\end{itemize}

\paragraph*{Výhody DevOps:}
\begin{itemize}
    \item \textbf{Rychlost:} Častější vydávání nových funkcí (místo jednou za půl roku klidně několikrát denně).
    \item \textbf{Spolehlivost:} Automatizované testy zachytí chyby dříve, než se dostanou k uživateli.
    \item \textbf{Škálovatelnost:} Díky cloudu a automatizaci lze snadno navyšovat výkon podle aktuální zátěže.
\end{itemize}

\paragraph*{Využití do Cloudu:}
\begin{itemize}
    \item \textbf{IaaS (Infrastructure as a Service):} Poskytovatel nabízí základní výpočetní zdroje – virtuální servery, úložiště a sítě. Uživatel si sám spravuje operační systém, middleware a aplikace (např. AWS EC2, Microsoft Azure VM).
    
    \item \textbf{PaaS (Platform as a Service):} Poskytovatel dodává kompletní platformu pro vývoj a běh aplikací. Vývojář se stará pouze o kód a data, nikoliv o správu OS nebo aktualizace runtime prostředí (např. Heroku, Google App Engine).

    \item \textbf{SaaS (Software as a Service):} Koncový software je poskytován přímo přes webový prohlížeč. Uživatel neřeší žádnou technickou stránku, pouze službu využívá (např. Office 365, Salesforce, Gmail).

    \item \textbf{FaaS (Function as a Service / Serverless):} Vývojář nasazuje pouze jednotlivé funkce (kousky kódu), které se spustí v reakci na určitou událost. Platí se pouze za čas běhu dané funkce (např. AWS Lambda, Azure Functions).

    \item \textbf{Modely nasazení (Deployment Models):}
    \begin{itemize}
        \item \textit{Veřejný cloud:} Sdílená infrastruktura u velkých poskytovatelů (nákladově efektivní).
        \item \textit{Privátní cloud:} Vyhrazená infrastruktura pro jednu organizaci (vyšší bezpečnost a kontrola).
        \item \textit{Hybridní cloud:} Kombinace obou přístupů – citlivá data v privátním cloudu, zbytek ve veřejném.
        \item \textit{Multi-cloud:} Využívání služeb od více poskytovatelů najednou (např. AWS i Azure), aby se předešlo závislosti na jednom dodavateli (vendor lock-in).
    \end{itemize}
\end{itemize}

\section{Webové služby a mikroslužby}
\begin{quote}
    \textbf{význam SOA, WS a mikroslužeb. Principy SOAP vs REST, XML a JSON. Nástroje na tvorbu mikroslužeb}
\end{quote}

\subsection{SOA (Service-Oriented Architecture)}
Architektura orientovaná na služby (SOA) je přístup k návrhu podnikového softwaru, kde jsou systémy tvořeny nezávislými, znovupoužitelnými službami. Tyto služby poskytují specifickou business funkcionalitu (např. ověření bonity klienta, fakturace) a komunikují spolu přes síť pomocí standardizovaných protokolů. 

\paragraph*{Klíčové koncepty:}
\begin{description}
    \item[Služba (Service):] Samostatný logický celek, který funguje jako "černá skříňka". Okolní systémy neznají její vnitřní implementaci, znají pouze její veřejné rozhraní (kontrakt).
    \item[ESB (Enterprise Service Bus):] Centrální komunikační sběrnice. Služby v SOA typicky nekomunikují bod-bod (napřímo), ale všechny se připojují na centrální ESB. Sběrnice zajišťuje směrování zpráv, transformaci formátů (např. z XML starého systému do JSON moderního) a zabezpečení.
    \item[Volná vazba (Loose Coupling):] Služby jsou navrženy tak, aby na sobě byly co nejméně závislé. Změna v jedné službě by neměla vyžadovat změny ve službách ostatních.
\end{description}

\paragraph*{Výhody a nevýhody:}
\begin{itemize}
    \item \textbf{Výhody:} Vysoká míra znovupoužitelnosti napříč celým podnikem. Skvělé pro integraci historických (legacy) systémů s novými aplikacemi.
    \item \textbf{Nevýhody:} Sběrnice ESB se při velkém nárůstu služeb stává úzkým hrdlem (bottleneck) a centrálním bodem selhání (SPOF - Single Point Of Failure). Komunikace přes těžkopádné protokoly (tradičně SOAP a XML) generuje velkou režii.
\end{itemize}

\paragraph*{Rozdíl mezi SOA a Mikroslužbami:}
Ačkoliv oba přístupy dělí systém na menší části, mají odlišnou filozofii:
\begin{itemize}
    \item \textbf{Rozsah a cíl:} SOA se snaží integrovat různé systémy v rámci \textit{celého podniku} (Enterprise level). Mikroslužby se naopak používají pro rozdělení \textit{jedné konkrétní aplikace} (Application level).
    \item \textbf{Způsob komunikace:} SOA spoléhá na "chytrou sběrnici" (Smart ESB), která obsahuje spoustu směrovací logiky. Mikroslužby využívají "hloupé roury" (Dumb pipes – např. REST API nebo jednoduché fronty zpráv) a veškerou chytrou logiku si drží uvnitř samotné služby.
    \item \textbf{Datová vrstva:} Služby v SOA velmi často přistupují ke sdílené obří podnikové databázi. Správně navržená mikroslužba by měla mít vždy svou vlastní, izolovanou databázi.
\end{itemize}

\subsection{WS (Web Services)}
Webové služby představují standardizovaný způsob integrace aplikací nezávisle na jejich platformě a programovacím jazyce (např. systém v Javě komunikuje se systémem v .NET). Tradiční WS jsou postaveny na standardech organizace W3C a komunikují výhradně pomocí XML.

\newpage
\paragraph*{Základní stavební kameny:}
\begin{description}
    \item[SOAP (Simple Object Access Protocol):] Komunikační protokol. Definuje formát zprávy zvaný \textit{SOAP Envelope} (obálka), která obsahuje hlavičku pro metadata (bezpečnost, transakce) a tělo pro samotná byznysová data.
    \item[WSDL (Web Services Description Language):] Detailní a striktní popis rozhraní (kontrakt) ve formátu XML. Přesně definuje, jaké operace (metody) služba nabízí a jaké datové typy přijímá a vrací. 
    \item[UDDI (Universal Description, Discovery, and Integration):] Adresář, do kterého mohou být služby registrovány, aby je ostatní systémy mohly dynamicky vyhledat. (V dnešní praxi se již téměř nepoužívá, ale je součástí definice).
\end{description}

\paragraph*{Srovnání Web Services (SOAP) vs. REST API:}
\begin{itemize}
    \item \textbf{Formát dat:} WS používá striktně XML. REST může používat cokoliv, nejčastěji odlehčený JSON.
    \item \textbf{Přístup:} WS je orientováno na \textit{akce/operace} (např. metoda \texttt{ZalozUcet()}). REST je orientován na \textit{zdroje/entity} a využívá standardní HTTP metody (POST na entitu \texttt{/ucet}).
    \item \textbf{Kontrakt a typování:} WS má striktní smlouvu (WSDL). Kód se z WSDL generuje automaticky. REST je volnější (i když dnes využívá např. standard OpenAPI/Swagger pro dokumentaci).
    \item \textbf{Výkon a režie:} Zprávy ve WS jsou kvůli struktuře XML mnohem větší a náročnější na zpracování. REST je mnohem rychlejší a vhodnější pro web a mobilní zařízení.
\end{itemize}

\subsection{Mikroslužby}
\label{sub:mikrosluzby2}
\begin{quote}
    \textbf{Předchozí kapitola:} \hyperref[sub:mikrosluzby]{Mikroslužby}
\end{quote}
Každá služba zajišťuje jednu konkrétní funkcionalitu. Nejčastěji se využívá REST API ke komunikaci.
Usnadňují škálování, vývoj, nasazení. Jsou odolnější vůči výpadku a dovolují využití různých
technologií.

\subsection{SOAP vs REST}
Výhodou SOAP je robustnostnost a standardizace na druhé straně REST je jednoduchý a flexibilní.
SOAP se typicky používá v rámci SOA pro složité a formálně definované webové služby.
REST nepoužívá žádný specifický protokol, ale staví na principu http.

\begin{description}
    \item[SOAP] -- Simple Object Access Protocol
    \item[REST] -- Representational State Transfer
\end{description}

\begin{table}[h]
\centering
\begin{tabular}{|l|l|l|}
\hline
\textbf{Vlastnost} & \textbf{SOAP} & \textbf{REST} \\ \hline
\textit{Protokol} & Vlastní protokol (nad HTTP/SMTP) & Architektonický styl nad HTTP \\ \hline
\textit{Formát zpráv} & Pouze XML & JSON, XML, HTML, text \\ \hline
\textit{Stavovost} & Může být stavový i bezstavový & Striktně bezstavový (Stateless) \\ \hline
\textit{Bezpečnost} & WS-Security (vysoká, na úrovni zprávy) & HTTPS (na úrovni transportu) \\ \hline
\textit{Výkon} & Nižší (režie XML obálky) & Vysoký (minimalistický přenos) \\ \hline
\textit{Definice} & WSDL (striktní kontrakt) & OpenAPI/Swagger (dokumentace) \\ \hline
\end{tabular}
\end{table}

\paragraph*{Srovnání datových formátů (XML vs. JSON):}

\begin{description}
    \item[XML (eXtensible Markup Language):] 
    Jazyk založený na značkách (tagy). Je hierarchický a velmi ukecaný.
    \begin{itemize}
        \item \textbf{Silné stránky:} Podpora schémat (XSD) pro validaci dat, podpora jmenných prostorů (namespaces), samopopisné značky.
        \item \textbf{Slabé stránky:} Velký objem dat (nadbytečné tagy), složitější parsování pro prohlížeče a mobilní zařízení.
    \end{itemize}

    \item[JSON (JavaScript Object Notation):] 
    Odlehčený formát založený na párech klíč–hodnota.
    \begin{itemize}
        \item \textbf{Silné stránky:} Minimální režie (malá velikost souborů), extrémně rychlé parsování v JavaScriptu, snadná čitelnost pro člověka.
        \item \textbf{Slabé stránky:} Chybí nativní podpora pro metadata a složité datové typy (např. datum je často jen řetězec), méně robustní validace oproti XSD.
    \end{itemize}
\end{description}

\paragraph*{Kdy co zvolit?}
\begin{itemize}
    \item \textbf{SOAP + XML:} Bankovní systémy, platební brány, státní správa. Tam, kde je prioritou bezpečnost, transakčnost a kde se rozhraní nemění desítky let.
    \item \textbf{REST + JSON:} Veřejná API, mobilní aplikace, webové frontendy, mikroslužby. Tam, kde je prioritou rychlost vývoje, snadná konzumace dat a výkon.
\end{itemize}

\subsection{Nástroje na tvorbu mikroslužeb}

\paragraph*{1. Frameworky:}
\begin{itemize}
    \item \textbf{Spring Boot (Java):} Nejpoužívanější framework v korporátní sféře. Nabízí moduly pro snadné vytvoření REST API a integraci s cloudovými službami (Spring Cloud).
    \item \textbf{Node.js / NestJS (JavaScript/TypeScript):} Populární díky asynchronnímu I/O modelu, což je ideální pro vysoce zatížená komunikační rozhraní.
    \item \textbf{Go (Golang):} Jazyk od Google navržený přímo pro cloudové služby. Vyniká rychlostí, nízkou spotřebou paměti a skvělou podporou souběžnosti.
\end{itemize}

\paragraph*{2. Infrastruktura a orchestrace:}
Mikroslužby by bez automatizované správy nebylo možné provozovat v desítkách či stovkách instancí.
\begin{itemize}
    \item \textbf{Docker:} Každá mikroslužba má svůj vlastní kontejner, který obsahuje vše potřebné k běhu.
    \item \textbf{Kubernetes (K8s):} Stará se o to, aby mikroslužby běžely, automaticky je škáluje při zátěži a restartuje při chybě.
    \item \textbf{Service Mesh (např. Istio):} Nadstavba nad Kubernetes, která řeší složitou síťovou komunikaci mezi službami (šifrování provozu, řízení toků, sledování latence).
\end{itemize}

\paragraph*{3. Správa konfigurace a Service Discovery:}
V dynamickém prostředí se IP adresy služeb neustále mění, proto jsou potřeba nástroje pro jejich nalezení.
\begin{itemize}
    \item \textbf{Service Registry (např. Consul, Netflix Eureka):} Telefonní seznam služeb. Každá nově spuštěná služba se zde nahlásí, aby ostatní věděli, kde ji najít.
    \item \textbf{API Gateway (např. Kong, NGINX):} Jediný vstupní bod pro klienty zvenčí. Směruje požadavky na správné vnitřní mikroslužby, řeší autentizaci a omezování zátěže (rate limiting).
\end{itemize}

\paragraph*{4. Monitoring a distribuované trasování:}
Když nastane chyba v řetězci deseti služeb, je těžké ji najít bez specializovaných nástrojů.
\begin{itemize}
    \item \textbf{Prometheus \& Grafana:} Nástroje pro sběr a vizualizaci metrik (např. vytížení CPU, počet požadavků za sekundu).
    \item \textbf{ELK stack (ElasticSearch, Logstash, Kybana):} Sběr, transformace a vizualizace logů.
    \item \textbf{Jaeger / Zipkin:} Nástroje pro distribuované trasování. Umožňují sledovat cestu jednoho konkrétního požadavku napříč všemi mikroslužbami, kterými prošel.
\end{itemize}

\section{NoSQL databáze a jejich základní principy}
\label{sec:nosql}
\begin{quote}
    \textbf{škálování, replikace, možnosti distribuce dat, Map\&Reduce, sharding, perzistence dat, zabezpečení, CAP teorém, BASE model}
\end{quote}
\hyperref[sub:dv]{Zpět} (Kapitola 4.)

NoSQL (Not Only SQL) databáze vznikly jako reakce na limity tradičních relačních databází při zpracování obrovských objemů nestrukturovaných dat a potřebě horizontálního škálování.
Upřednostňují výkon dostupnost a flexibilitu před integritou a transakčností.
Zvládají rychlý zápis a čtení, které by klasické RDBMS nezvládalo.

\subsection{Škálování}
Škálování je schopnost systému zvládat rostoucí zátěž (počet uživatelů, objem dat) navýšením dostupných prostředků.

\paragraph*{1. Vertikální škálování (Scale-Up)}
Znamená navýšení výkonu stávajícího serveru - přidání více CPU, RAM nebo rychlejších disků.
\begin{itemize}
    \item \textbf{Výhody:} Jednoduchost (aplikace se nemusí měnit, stále běží na jednom stroji), žádná režie na síťovou komunikaci mezi uzly.
    \item \textbf{Nevýhody:} Existuje pevný hardwarový strop (limit základní desky), každé další navýšení výkonu je extrémně drahé a obvykle vyžaduje odstávku systému pro upgrade.
\end{itemize}

\paragraph*{2. Horizontální škálování (Scale-Out)}
Znamená přidávání dalších, levnějších strojů (uzlů) do clusteru. Zátěž se pak rozděluje mezi ně.
\begin{itemize}
    \item \textbf{Výhody:} Prakticky neomezený růst, vysoká dostupnost (když jeden uzel selže, ostatní pracují dál), možnost využít levnější hardware.
    \item \textbf{Nevýhody:} Vyšší komplexnost - systém musí být navržen jako bezstavový (stateless), je nutné řešit synchronizaci dat a použít \textbf{Load Balancer} (rozdělovač zátěže).
\end{itemize}

\subsection{Replikace a sharding}
Replikace zajišťuje redundanci dat, a tak existuje v případě ztráty nebo poškození
záložní kopie. Sharding je rozdělení dat na více uzlů, což je základní princip
horizontálního škálování.

\paragraph*{1. Replikace}
Proces udržování stejných dat na více serverech současně.
\begin{itemize}
    \item \textbf{Master-Slave (Leader-Follower):} Jeden uzel je hlavní (Master) a přijímá zápisy. Ostatní uzly (Slaves) kopírují jeho data a slouží pouze pro čtení. 
        \begin{itemize}
            \item \textit{Výhoda:} Extrémní zrychlení čtení (typické pro weby - hodně se čte, málo píše).
            \item \textit{Nevýhoda:} Zapisovat múže pouze Master.
            \item Např. Mongo
        \end{itemize}
    \item \textbf{Multi-Master:} Zápis je možný do libovolného uzlu. 
        \begin{itemize}
            \item \textit{Výhoda:} Všechny uzly mohou číst i zapisovat.
            \item \textit{Nevýhoda:} Velmi složité řešení konfliktů (co když dva lidé změní stejný řádek na různých serverech ve stejnou vteřinu).
            \item Např. Cassandra
        \end{itemize}
    \item \textbf{Hlavní přínos:} Vysoká dostupnost (Availability) - když jeden server shoří, data máme jinde.
    \item \textbf{Synchronní/asynchronní:} Synchronní zápis je potvrzen až po zapsání do všech uzlů. Asynchronní zápis je potvrzen ihned, replikace proběhne později.
\end{itemize}

\paragraph*{2. Sharding (Horizontální dělení)}
Rozdělení jedné obří tabulky na několik menších částí (shardy), které jsou fyzicky uloženy na různých serverech. Žádný server pak nemá kompletní kopii všech dat.
\begin{itemize}
    \item \textbf{Shard Key:} Klíč (např. ID zákazníka nebo region), podle kterého se určuje, na kterém serveru data leží. Např. uživatelé z Evropy jsou na serveru A, uživatelé z USA na serveru B.
    \item \textbf{Výhody:} Umožňuje lineární růst databáze. Zápis se rozprostře mezi více strojů, takže neexistuje jeden bottleneck.
    \item \textbf{Nevýhody:} Velká komplexnost. Dotazy přes více shardů (JOIN) jsou velmi pomalé a náročné na režii.
\end{itemize}

\paragraph*{Možnosti distribuce dat:}
Označuje způsob jakým jsou data uložena na více fyzických nebo logických uzlech.
\begin{itemize}
    \item \textbf{HashBased} -- Data jsou distribuována podle výstupu hashovací funkce nad klíčem což zajišťuje rovnoměrné rozložení.
    \item \textbf{RangeBased} -- Data jsou rozdělena podle hodnoty klíče do určitých rozsahů.
    \item \textbf{Geografická distribuce} Každý uzel je v jiném datacentru nebo geografické oblasti. Ukládá data ze své oblasti.
\end{itemize}

\subsection{Map\&Reduce}
MapReduce je programovací model pro distribuované zpracování obrovských objemů dat (BigData) na velkých clusterech. Tento koncept zpopularizoval Google a je základem technologií jako Apache Hadoop.

\paragraph*{Princip fungování:}
Celý proces je rozdělen do dvou hlavních fází, mezi kterými probíhá automatické míchání dat:

\begin{description}
    \item[1. Map (Mapování):] Vstupní data jsou rozdělena na menší kusy (případně již jsou uložena na různých uzlech) a každý uzel zpracuje svou část a vygeneruje dvojice \textbf{klíč–hodnota} (key-value pairs).
    \item[2. Shuffle and Sort (Míchání):] Systém automaticky seskupí všechny hodnoty se stejným klíčem k sobě a pošle je na stejný uzel pro závěrečné zpracování.
    \item[3. Reduce (Redukce):] Uzly přijmou seskupená data pro daný klíč a provedou agregační operaci (např. součet, průměr, výběr maxima), čímž vznikne finální výsledek.
\end{description}

\paragraph*{Příklad (Počítání slov - Word Count):}
Máme miliony dokumentů a chceme zjistit počet výskytů každého slova:
\begin{itemize}
    \item \textbf{Map:} Každý stroj čte text a pro každé slovo "ahoj" vytvoří dvojici \texttt{(ahoj, 1)}. (číslo je počet nálezů slova)
    \item \textbf{Shuffle:} Všechny dvojice \texttt{(ahoj, *)} se "sletí" na jeden konkrétní stroj.
    \item \textbf{Reduce:} Tento stroj sečte všechny jedničky u klíče "ahoj" a vypíše výsledek \texttt{(ahoj, 1548)}.
\end{itemize}

\paragraph*{Výhody a technologie:}
\begin{itemize}
    \item \textbf{Škálovatelnost:} Stačí přidat další servery a výkon roste lineárně.
    \item \textbf{Odolnost proti chybám:} Pokud jeden uzel během výpočtu selže, systém jeho práci automaticky přidělí jinému stroji.
\end{itemize}

\subsection{Perzistence dat a zabezpečení}


\paragraph*{Perzistence dat a distribuovaný přístup:}
Zatímco relační databáze striktně dodržují pravidla ACID (Atomicita, Konzistence, Izolace, Trvalost), distribuované NoSQL databáze se často řídí \textbf{teorémem CAP} (Konzistence, Dostupnost, Odolnost proti rozdělení sítě) a principy \textbf{BASE} (Basically Available, Soft state, Eventual consistency). To znamená, že systém raději upřednostní dostupnost a rychlost před okamžitou konzistencí napříč všemi uzly.
Mechanismy perzistence v NoSQL lze rozdělit do několika přístupů:
\begin{itemize}
    \item \textbf{In-memory s asynchronním zápisem:} Databáze (např. Redis) drží data primárně v operační paměti (RAM) pro maximální rychlost. Perzistence je řešena občasným vytvořením snímku (Snapshotting) na pevný disk nebo průběžným logováním každé operace (Append-Only File).
    \item \textbf{Write-Ahead Logging (WAL):} Předtím, než se data zapíšou do samotné datové struktury, jsou nejprve sekvenčně zaznamenána do transakčního logu na disku. Tím je zaručeno, že v případě výpadku proudu se data neztratí a systém je z logu obnoví (využívá např. Cassandra nebo MongoDB).
\end{itemize}

\paragraph*{Zabezpečení:}
Historicky byly NoSQL databáze navrhovány s primárním důrazem na výkon a škálovatelnost. Předpokládalo se, že poběží v uzavřených a důvěryhodných vnitřních sítích, a proto v počátcích často postrádaly robustní zabezpečení (např. výchozí nastavení nevyžadovalo heslo). Moderní NoSQL systémy však již implementují komplexní bezpečnostní vrstvy:

\begin{itemize}
    \item \textbf{Autentizace a Authorizace:} Moderní systémy využívají standardní mechanismy pro ověření identity (SCRAM, integrace s LDAP nebo Kerberos). Pro řízení přístupu se standardně používá \textbf{RBAC} (Role-Based Access Control) -- uživatelům nebo aplikacím jsou přidělovány role s přesně definovanými oprávněními (např. pouze čtení z konkrétní kolekce).
    \item \textbf{Šifrování (Encryption):}
    \begin{itemize}
        \item \textbf{In-transit (při přenosu):} Komunikace mezi klientem a databází, stejně jako mezi samotnými uzly v clusteru, je šifrována pomocí TLS/SSL protokolů, aby se zabránilo odposlechu (man-in-the-middle útokům).
        \item \textbf{At-rest (v úložišti):} Data fyzicky uložená na pevných discích jsou šifrována (např. pomocí standardu AES-256). K datům nelze přistoupit pouhým zkopírováním databázových souborů bez šifrovacího klíče.
    \end{itemize}
    \item \textbf{Síťová izolace a Auditing:} Databázové uzly jsou chráněny firewally a přístup k nim je povolen pouze z aplikačních serverů (nikoliv z veřejného internetu). Většina enterprise verzí NoSQL databází navíc podporuje podrobné auditování, které zaznamenává, kdo, kdy a k jakým datům přistupoval.
\end{itemize}

\subsection{CAP teorém a BASE model}
Tradiční relační databáze se spoléhají na vlastnosti ACID (Atomicita, Konzistence, Izolace, Trvalost), které zaručují absolutní spolehlivost transakcí. V prostředí obrovských distribuovaných systémů (BigData) je však striktní dodržování ACID velmi náročné na výkon a často nemožné. Proto NoSQL databáze využívají odlišné teoretické modely -- CAP teorém a principy BASE.

\paragraph*{CAP teorém (Brewerův teorém):}
Tento teorém definuje tři základní vlastnosti distribuovaných systémů. Podle něj může distribuované datové úložiště současně garantovat \textbf{nejvýše dvě} z těchto tří vlastností:

\begin{description}
    \item[C -- Consistency (Konzistence):] Každý požadavek na čtení vrátí buď nejnovější zapsaná data, nebo chybu. Všichni uživatelé vidí v daný okamžik naprosto stejná data, nezávisle na tom, ke kterému uzlu (serveru) jsou připojeni.
    \item[A -- Availability (Dostupnost):] Každý požadavek, který systém přijme, dostane platnou odpověď (bez chybové hlášky). Není ale zaručeno, že tato odpověď obsahuje nejnovější verzi dat.
    \item[P -- Partition Tolerance (Odolnost proti rozdělení sítě):] Systém nadále funguje i v případě, že dojde k výpadku komunikace (rozdělení sítě) mezi jednotlivými uzly.
\end{description}

\textbf{Volba v praxi:} Ačkoliv se často říká, že si \uv{vybíráme dvě ze tří}, v reálném světě jsou výpadky sítě (P) nevyhnutelné. Distribuovaný systém proto \textit{musí} být odolný proti rozdělení (Partition Tolerant). Skutečná volba při návrhu NoSQL databáze tedy spočívá v tom, co systém udělá při výpadku sítě:
\begin{itemize}
    \item \textbf{CP systémy (Konzistence + Odolnost):} Při výpadku sítě systém raději odmítne odpovědět (obětuje dostupnost), aby nevrátil zastaralá data (např. MongoDB, HBase).
    \item \textbf{AP systémy (Dostupnost + Odolnost):} Systém odpoví vždy, i za cenu, že klientům pošle starší data, protože uzly se nestihly synchronizovat (např. Cassandra, CouchDB).
\end{itemize}

\paragraph*{Model BASE:}
Zatímco ACID je model pro relační databáze, NoSQL databáze (zejména ty typu AP) se řídí modelem BASE. Jde o akronym tří konceptů, které upřednostňují dostupnost před okamžitou konzistencí:

\begin{description}
    \item[BA -- Basically Available (Základní dostupnost):] Systém garantuje vysokou dostupnost dat. I když selže jeden uzel nebo část sítě, databáze jako celek stále odpovídá na požadavky uživatelů.
    \item[S -- Soft state (Měkký/proměnlivý stav):] Vzhledem k chybějící garanci okamžité konzistence se stav systému může měnit v čase, a to i v případě, že nepřicházejí žádné nové požadavky na zápis (uzly se totiž synchronizují asynchronně na pozadí).
    \item[E -- Eventual consistency (Konečná konzistence):] Systém negarantuje, že data budou konzistentní okamžitě po zápisu. Garantuje však, že pokud do systému přestanou přicházet nové zápisy, po uplynutí určité doby se všechny uzly zkoordinují a data budou napříč systémem konzistentní.
\end{description}

\section{Data v distribuovaných systémech} 
\begin{quote}
    \textbf{BigData a jejich zdroje, dataset a jeho životní cyklus, otevřená data, trend velkých datasetů, 5V, typy BigData a datové formáty, problémy BigData}
\end{quote}

\subsection{BigData}
Velké, rychle vznikající různorodé objemy dat, které není možné efektivně zpracovávat tradičními metodami a nástroji.
Typicky několik TB až PB. Ve skutečnosti neexistuje pevná definice toho, co jsou
BigData, jedná se o buzzword, pojem se stále vyvíjí, rychle se mění.

\begin{quote}
    \textbf{Gartner:} „BigData“ jsou informační zdroje charakterizované velkým objemem, vysokou rychlostí a/nebo velkou rozmanitostí,
které vyžadují nové způsoby zpracování, aby umožnily lepší rozhodování, získávání poznatků a optimalizaci procesů.
\end{quote}

\subsubsection{Zdroje:}
\begin{itemize}
    \item Senzory a IoT zařízení
    \item Web a mobilní aplikace -- logy, pohyby myší, interakce, data o chování 
    uživatelů
    \item Sociální sítě -- texty, obrázky, videa
    \item Finanční a obchodní transakce -- platby, objednávky, interakce v eshopu
    \item Veřejné zdroje a otevřená data
    \item Zdravotnictví a biomedicína
    \item Multimediální obsah
\end{itemize}

\subsection{Dataset}

Dataset (datová sada) je strukturovaná kolekce souvisejících dat, obvykle uspořádaná do tabulkové podoby (kde řádky reprezentují jednotlivé záznamy a sloupce jejich atributy), případně do podoby textových korpusů, grafů či složek s multimédii. Technicky jsou datasety uloženy v různých formátech, od jednoduchých souborů (CSV, JSON) až po komplexní relační či NoSQL databáze.

\subsubsection*{Životní cyklus datasetu}
Práce s daty se řídí procesem, který pokrývá jejich existenci od vzniku po zánik. Hlavní fáze zahrnují:
\begin{itemize}
    \item \textbf{Sběr (monitoring, scraping):} Získávání surových dat z různých zdrojů (uživatelské vstupy, senzory, web scraping, logy aplikací).
    \item \textbf{Organizace:} Organizace dat do datasetu. Řeší se například jmenné konvence, čištění, předspracování.
    \item \textbf{Ukládání:} Zabezpečené uložení dat do databází nebo jiným způsobem.
    \item \textbf{Prohledávání:} Průzkum dat například dotazy nad databází.
    \item \textbf{Sdílení:} Sdílení dat ostatním, např. registry otevřených dat.
    \item \textbf{Analýza a využití:} Vytěžování informací (Data Mining), podpora rozhodování (BI) nebo trénování modelů strojového učení (ML). Např. Grafana, Dataprep.
    \item \textbf{Vizualizace:} Prezentace výsledků získaných z dat. Např. dashboardy (Datastudio).
    \item \textbf{Archivace či zničení:} Odložení historických dat z důvodu úspory drahého úložiště nebo bezpečné smazání dat na základě legislativních požadavků (např. expirace souhlasu dle GDPR).
\end{itemize}

\subsection{Otevřená data (Open Data)}
Otevřená data jsou informace a datasety bezplatně zpřístupněné veřejnosti k volnému využití (bez licenčních či autorských omezení). Aby data byla skutečně ,,otevřená'', musí být publikována ve \textbf{strojově čitelném formátu} (např. CSV nebo XML, nikoliv skenovaný obrázek v PDF) a nesmí být vázána na proprietární software.
Cílem otevřených dat je podpora transparentnosti (např. veřejné rozpočty, jízdní řády v Národním katalogu otevřených dat v ČR) a podpora inovací, neboť na nich mohou nezávislí vývojáři stavět nové služby a aplikace.
\href{https://5stardata.info/cs/}{Web} na vysvětlení úrovní.

\subsection{Trend velkých datasetů (BigData)}
V současnosti čelíme exponenciálnímu nárůstu produkce dat. Zpracování velkých 
datasetů vyžaduje využití cloudu a distribuovaných výpočetních systémů 
(např. Apache Hadoop, Spark). Obrovské datasety jsou v současnosti naprosto 
klíčové pro trénování modelů umělé inteligence, což však přináší i nové výzvy v 
oblasti kvality dat, biasu (zaujatosti) a ochrany soukromí. Data se také stala komoditou.
\begin{itemize}
\item Spojování malých datasetů do velkých.
\item Vyšší nároky na výkon DB.
\item Objevení i zdánlivě neexistujících korelací.
\item Killer-feature pro business intelligence!
\item Klíčovými aspekty platformy BigData jsou integrace, analýzy, vizualizace a zabezpečená správa.
\end{itemize}

\subsection{Koncept 5V}
Původní charakteristika velkých dat (Volume, Velocity, Variety) se dnes standardně rozšiřuje o další dvě dimenze (Veracity, Value), které definují úspěšnost práce s BigData:
\begin{itemize}
    \item \textbf{Volume (Objem):} Obrovské množství dat (terabyty až petabyty), které již nelze zpracovat na běžných počítačích či v tradičních databázích.
    \item \textbf{Velocity (Rychlost):} Nutnost data sbírat a analyticky zpracovávat v reálném čase (např. IoT senzory, detekce bankovních podvodů).
    \item \textbf{Variety (Rozmanitost):} Drtivá většina moderních dat je nestrukturovaná (obrázky, videa, příspěvky ze sociálních sítí, audio nahrávky).
    \item \textbf{Veracity (Věrohodnost/Pravdivost):} Týká se kvality a spolehlivosti dat. Při obrovských objemech data často obsahují šum, nekonzistence, duplicity nebo chyby. Pokud je špatný vstup, je špatný i výstup analýzy. Popisuje reálnou shodu s fakty a přesností.
    \item \textbf{Value (Hodnota):} Data sama o sobě nemají žádnou cenu. Hodnota představuje schopnost vytěžit z masivních datasetů užitečné znalosti, které přinesou byznysovou výhodu nebo zefektivní procesy.
\end{itemize}

\subsection{Typy BigData a datové formáty}
Z hlediska formálního uspořádání dělíme data do tří základních kategorií:
\begin{itemize}
    \item \textbf{Strukturovaná data:} Mají přesně definovaný datový model a schéma. Lze je snadno ukládat a dotazovat pomocí SQL. \textit{Formáty: CSV, tabulky v relačních databázích (MySQL, PostgreSQL).}
    \item \textbf{Polostrukturovaná data:} Nemají rigidní tabulkovou strukturu, ale obsahují hierarchii, klíče nebo značky (tagy) pro oddělení datových prvků. \textit{Formáty: JSON, XML, YAML (typické pro webová API a NoSQL databáze).}
    \item \textbf{Nestrukturovaná data:} Tvoří zhruba 80~\% všech moderních dat. Nemají předem definovaný model a nelze je snadno tabulkově prohledávat. Analýza často vyžaduje nástroje umělé inteligence. \textit{Příklady: textové dokumenty (PDF), e-maily, obrázky (JPEG), video (MP4), audio stopy.}
\end{itemize}

\subsection{Problémy BigData}
\begin{itemize}
    \item \textbf{Ukládání a škálování:} Není možné využívat centrální servery. Nutnost horizontálního škálování (distribuované systémy a cloudy), což razantně zvyšuje složitost architektury. Pro dotazování v databázi: cacheování, indexace
    \item \textbf{Integrace dat:} Obtížné propojování dat (Data Integration) z izolovaných systémů (tzv. datová sila), které navíc využívají zcela odlišné formáty a struktury.
    \item \textbf{Bezpečnost a ochrana soukromí:} Velká úložiště jsou lákavým cílem kyberútoků. Sběr osobních dat naráží na přísnou legislativu (GDPR), což vyžaduje spolehlivou anonymizaci či pseudonymizaci.
    \item \textbf{Výkon a latence:} IoT zařízení mají nízký výkon. Neustálý přenos dat po síti může síť přetížit. Nutná optimalizace přenosu.
\end{itemize}

\newpage

\section{Struktura a architektura systému UNIX / GNU / Linux}
\begin{quote}
    \textbf{struktura systému, systém souborů, systém procesů, služby OS, shelly, příkazy shellu, popis jádra, datové struktury jádra, systém vyrovnávací paměti, subsystém V/V, správa paměti, operační systémy reálného času (základní charakteristika, hlavní faktury, definice, hard a soft RTOS, RMS, EDF, příklady RTOS)}
\end{quote}
\begin{itemize}
    \item \textbf{UNIX:} Komerční operační systém vytvořený v 70. letech v Bellových laboratořích (AT\&T). Přinesl revoluční myšlenky: víceúlohovost (multitasking), víceuživatelský přístup a filozofii ,,vše je soubor''. Původní kód UNIXu je uzavřený (proprietární) a dnes z něj vycházejí např. systémy macOS nebo FreeBSD.
    \item \textbf{GNU (GNU's Not Unix):} Projekt založený v roce 1983 (Richard Stallman) s cílem vytvořit svobodný (open-source) operační systém, který by se choval jako UNIX, ale neobsahoval jeho komerční kód. Projekt vytvořil klíčové uživatelské nástroje (příkazový řádek bash, překladač GCC, základní utility), ale chybělo mu funkční jádro.
    \item \textbf{Linux:} Jádro (kernel) operačního systému, které v roce 1991 napsal Linus Torvalds.
    Spojením jádra \textbf{Linux} a uživatelských nástrojů \textbf{GNU} vznikl plně funkční a svobodný operační systém, který se dnes zkráceně, označuje jako ,,Linux'' (správně by mělo být GNU/Linux).
\end{itemize}


\subsection{Struktura systému}
Linux využívá monolitické jádro a striktně odděluje práva hardwaru a běžných uživatelů. Architektura se dělí do několika vrstev:
\begin{itemize}
    \item \textbf{Hardware:} Fyzické komponenty (CPU, RAM, disky, síťové karty).
    \item \textbf{Jádro (Kernel Space):} Běží v privilegovaném režimu. Řídí přístup k hardwaru, spravuje paměť, plánuje procesorový čas pro aplikace a obsahuje ovladače zařízení. Jádro je program, který je zaveden po zapnutí počítače do operační paměti.
    \item \textbf{Systémová volání (System Calls):} Rozhraní (API) mezi jádrem a uživatelským prostorem. Pokud chce běžná aplikace zapsat soubor na disk, nesmí sáhnout přímo na HW, ale musí o to požádat jádro právě přes systémové volání.
    \item \textbf{Uživatelský prostor (User Space):} Neprivilegovaný režim, ve kterém běží knihovny (např. glibc), příkazový interpret (Shell), služby na pozadí (démoni), grafické prostředí (X11/Wayland) a samotné uživatelské aplikace.
\end{itemize}

\subsection{Systém souborů (File System)}
Základním principem je, že \textbf{,,vše je soubor''} - jako soubor se v Linuxu tváří textové dokumenty, složky, ale i hardwarová zařízení (např. pevný disk jako \texttt{/dev/sda}).
\begin{itemize}
    \item \textbf{Hierarchie:} Vše vychází z jednoho kořenového adresáře označovaného lomítkem (\texttt{/}), další disky se připojují (mountují) do prázdných složek v rámci této jedné hierarchie.
    \item \textbf{I-uzly (I-node):} Datová struktura, která nese metadata o souboru (vlastník, velikost, přístupová práva, časy modifikace, typ souboru a hlavně ukazatele na fyzické bloky na disku, kde jsou data skutečně uložena). I-uzel neobsahuje název souboru - ten je uložen pouze v adresáři.
    \item \textbf{Odkazy (Linky):}
    \begin{itemize}
        \item \textit{Tvrdý odkaz (Hard link):} Více názvů souborů ukazuje na stejný i-uzel. Data se smažou z disku až po smazání posledního tvrdého odkazu.
        \item \textit{Symbolický odkaz (Symlink):} Obdoba zástupce ve Windows. Soubor, který obsahuje pouze cestu k jinému souboru.
    \end{itemize}
    \item \textbf{Oprávnění:} Každý soubor má vlastníka (User) a skupinu (Group). Práva se dělí na čtení (\texttt{r}), zápis (\texttt{w}) a spouštění (\texttt{x}) pro vlastníka, skupinu a všechny ostatní (Others).
\end{itemize}
\subsubsection*{Základní typy souborů}
\label{prg:typysoub}
\begin{itemize}
    \item \textbf{Běžný soubor} (značí se pomlčkou \texttt{-}): Nejobvyklejší typ. Patří sem textové dokumenty, zdrojové kódy, obrázky, komprimované archivy i spustitelné binární programy.
    \item \textbf{Adresář} (značí se \texttt{d} - Directory): Z pohledu jádra jde o \textit{speciální soubor}, který neobsahuje běžná data, ale pouze tabulku. V této tabulce jsou zapsány názvy souborů (které se v adresáři nacházejí) a k nim přiřazená čísla i-uzlů.
    \item \textbf{Symbolický odkaz} (značí se \texttt{l} - Link): Soubor obsahující pouze textovou cestu k jinému cílovému souboru (obdoba zástupce).
    \item \textbf{Soubory zařízení (Device files):} Slouží k přímé komunikaci s hardwarem. Nacházejí se typicky ve složce \texttt{/dev}. Dělí se na dva druhy:
    \begin{itemize}
        \item \textit{Znaková zařízení (značí se \texttt{c} - Character):} Přenáší data znak po znaku (proud bajtů) bez vyrovnávací paměti. Např. klávesnice, myš, terminál.
        \item \textit{Bloková zařízení (značí se \texttt{b} - Block):} Přenáší data ve větších blocích (např. 4 KB) přes vyrovnávací paměť. Např. pevné disky (\texttt{/dev/sda}) nebo USB flash disky.
    \end{itemize}
    \item \textbf{Pojmenovaná roura} (značí se \texttt{p} - Pipe / FIFO): Speciální soubor sloužící k jednosměrné komunikaci mezi dvěma nezávislými procesy.
    \item \textbf{Soket} (značí se \texttt{s} - Socket): Speciální soubor umožňující obousměrnou pokročilou komunikaci mezi procesy, a to lokálně nebo přes síť.
\end{itemize}
Zpět na \hyperref[prg:vv]{subsystém V/V}.

\subsection{Systém procesů}
Proces je běžící instance programu. Linux je plně víceúlohový, takže procesor přepíná mezi stovkami procesů (Context switching).
Všechny porcesy jsou zapsány v tabulce procesů.
\begin{itemize}
    \item \textbf{Atributy procesu:} Každý proces má unikátní identifikátor \textbf{PID} (Process ID), identifikátor rodiče \textbf{PPID} (Parent PID) a práva uživatele (\textbf{UID} - uživatel, \textbf{GID} - skupina), který jej spustil.
    \item \textbf{Vznik procesu:} Vytvoření nového procesu probíhá ve dvou krocích. Voláním \texttt{fork()} vytvoří rodičovský proces svou přesnou kopii (potomka). Potomek následně použije volání \texttt{exec()}, kterým svůj obsah nahradí novým programem.
    \item \textbf{Init proces (systemd):} První proces spuštěný při startu jádra. Má vždy \textbf{PID 1} a je přímým nebo nepřímým rodičem všech ostatních procesů v systému.
    \item \textbf{Meziprocesová komunikace (IPC):} Procesy si mohou vyměňovat data nebo se synchronizovat. Využívají se k tomu \textit{roury} (Pipes - znak \texttt{|}, přesměrování výstupu jednoho programu na vstup dalšího), signály (např. \texttt{SIGKILL} pro okamžité ukončení procesu) nebo sdílená paměť.
    \item \textbf{/proc souborový systém:} Pseudo-souborový systém poskytující rozhraní k interním údajům o běžících procesech a systému. (\texttt{/proc/cpuinfo} - informace o CPU, \texttt{/proc/meminfo}, \texttt{/proc/uptime})
\end{itemize}
\subsubsection*{Stavy procesů}
\label{prg:stavy}
\begin{itemize}
    \item Běžící nebo připravený běžet
    \item Spící připravený na probuzení signálem
    \item Čeká bez možnosti přerušení signálem
    \item Zastavený
    \item Zombie
\end{itemize}
\subsubsection*{Datová struktura procesu (\texttt{task\_struct})}
\label{prg:taskstruct}
V jádře Linuxu je každý proces (i vlákno) interně reprezentován  strukturou zvanou \textbf{\texttt{task\_struct}}. Jádro si tyto struktury udržuje v paměti a obsahuje v nich naprosto vše, co o procesu potřebuje vědět:

\begin{itemize}
    \item \textbf{Stav procesu:} \hyperref[prg:stavy]{Stavy procesu} 
    \item \textbf{Identifikátory:} \texttt{pid} (číslo procesu), \texttt{tgid} (číslo skupiny vláken) a údaje o vlastníkovi (UID/GID) pro ověření oprávnění.
    \item \textbf{Rodokmen:} Ukazatele na rodičovský proces (\texttt{parent}) a seznam potomků (\texttt{children}).
    \item \textbf{Správa paměti (\texttt{mm\_struct}):} Odkaz na datovou strukturu definující virtuální adresní prostor procesu (kde má proces kód, data a zásobník).
    \item \textbf{Otevřené soubory (\texttt{files\_struct}):} Tabulka tzv. \textit{file deskriptorů} (ukazatele na konkrétní soubory, roury či síťové sokety, které má proces zrovna otevřené).
    \item \textbf{Plánování a kontext CPU:} Priorita procesu a paměťový prostor, kam jádro uloží stavy registrů procesoru při přepínání úloh (Context Switch). Číselné statistiky např. kolik času strávil v uživatelském / jadru. Také \texttt{thread\_info} s informacemi o kontextu při přepínání vláken.
\end{itemize}
Zpět na \hyperref[prg:datastructs]{Datové struktury jádra}.

\subsection{Služby OS}
Služby operačního systému jsou základní softwarové funkce, které zajišťují chod počítače, správu hardwaru a poskytují prostředí pro spouštění aplikací.
\begin{itemize}
    \item Správa procesů
    \item Správa paměti
    \item Správa souborového systému
    \item Správa zařízení
    \item Správa uživatelů a práv
    \item Síťové služby
    \item Synchronizace a komunikace mezi procesy
    \item Cachování diskového přístupu
    \item Virtuální paměť se stránkováním do tzv. odkládacího prostoru swap space
    \item Ukládání obsahu paměti na disk core dump
    \item Pravý hardwarový paralelismus
    \item Stránkování na vyžádání
\end{itemize}

\subsection{Shelly}

\textbf{Shell} (příkazový interpret, skutečný programovací jazyk) je textové uživatelské rozhraní k OS. Čte příkazy zadané uživatelem (nebo ze skriptu), interpretuje je a spouští příslušné programy.

\subsubsection*{Běžné typy shellů}
\begin{itemize}
    \item \textbf{sh (Bourne shell):} Historický základ, původní shell systému UNIX.
    \item \textbf{bash (Bourne Again shell):} Výchozí a nejrozšířenější shell ve většině linuxových distribucí.
    \item \textbf{zsh, fish, csh, a další:} Moderní shelly s pokročilým našeptáváním, obarvováním syntaxe a bohatými možnostmi přizpůsobení.
\end{itemize}

\subsubsection*{Typy příkazů}
\begin{itemize}
    \item \textbf{Interní (vestavěné) příkazy:} Jsou přímou součástí samotného shellu. Při jejich spuštění nevzniká nový proces. Příklad: \texttt{cd} (změna adresáře), \texttt{echo}, \texttt{export} (nastavení proměnné).
    \item \textbf{Externí příkazy:} Jsou to samostatné binární soubory uložené na disku. Shell je hledá v adresářích definovaných v systémové proměnné \textbf{\texttt{\$PATH}}. Při spuštění shell vytvoří nový proces (přes \texttt{fork} a \texttt{exec}). Příklad: \texttt{ls}, \texttt{grep}, \texttt{cat}.
\end{itemize}

\subsubsection*{Standardní proudy a přesměrování}
Každý spuštěný proces v Linuxu má automaticky otevřené tři datové proudy (soubory): \textbf{stdin} (0 - standardní vstup, typicky klávesnice), \textbf{stdout} (1 - standardní výstup, obrazovka) a \textbf{stderr} (2 - chybový výstup, obrazovka). Shell s nimi umí manipulovat:
\begin{itemize}
    \item \textbf{Přesměrování výstupu:} Znak \texttt{>} (uloží výstup do souboru, původní obsah přepíše), znaky \texttt{$>>$} (připojí výstup na konec souboru).
    \item \textbf{Roury (Pipes):} Znak \texttt{|}. Propojuje procesy – standardní výstup prvního příkazu se stává standardním vstupem druhého (např. \texttt{ls -l | grep "txt"} vypíše jen soubory obsahující "txt").
\end{itemize}

\subsubsection*{Přehled nejpoužívanějších příkazů}
\begin{itemize}
    \item \textbf{Práce se soubory/adresáři:} \texttt{ls} (výpis), \texttt{cd} (změna složky), \texttt{pwd} (aktuální cesta), \texttt{cp} (kopírování), \texttt{mv} (přesun/přejmenování), \texttt{rm} (mazání).
    \item \textbf{Zpracování textu:} \texttt{cat} (výpis obsahu souboru), \texttt{grep} (filtrování textu dle vzoru/regexu), \texttt{less} (stránkování textu).
    \item \textbf{Správa procesů:} \texttt{ps} (výpis běžících procesů), \texttt{top} / \texttt{htop} (interaktivní sledování zátěže), \texttt{kill} (zaslání signálu procesu, typicky pro ukončení).
    \item \textbf{Oprávnění:} \texttt{chmod} (změna práv r/w/x), \texttt{chown} (změna vlastníka).
\end{itemize}

\subsection{Jádro OS}
Jádro (neboli kernel) je programové vybavení, které pracuje
na počítači bez jakékoliv další programové podpory, představuje množinu obslužných programů trvale
umístěných v paměti, k nimž má každý proces přístup,
základem operačního systému a odpovídá za správu všech hardwarových prostředků počítače, jako jsou disky, síťová
rozhraní atd.
Je velkým rozhraním mezi uživatelem a technickým
vybavením počítače.
Udržuje a podporuje v činnosti systém souborů a běžící
procesy.
\subsubsection*{Popis jádra}
\begin{itemize}
    \item \textbf{Monolitická architektura:} Celé jádro (včetně většiny ovladačů a souborových systémů) běží v jednom velkém paměťovém bloku. To zajišťuje maximální výkon.
    \item \textbf{LKM (Loadable Kernel Modules):} Přestože je Linux monolitický, je \textit{modulární}. Ovladače hardwaru lze do běžícího jádra dynamicky nahrávat a odebírat za běhu (příkazy \texttt{modprobe}, \texttt{lsmod}), aniž by bylo nutné systém restartovat.
\end{itemize}

\subsubsection*{Rozdíl mezi monolitickým jádrem a mikrojádrem}
Základní rozdíl spočívá v tom, kolik služeb běží v privilegovaném režimu (Kernel Space).

\begin{itemize}
    \item \textbf{Monolitické jádro (např. Linux):} Všechny služby OS (správa paměti, souborové systémy, ovladače HW) tvoří jeden velký celek a běží v kernel space.
    \begin{itemize}
        \item \textit{Výhoda:} Maximální výkon a rychlost (komunikace probíhá přímo v paměti bez režie).
        \item \textit{Nevýhoda:} Nižší stabilita (chyba nebo zásek v jediném ovladači může způsobit pád celého OS - tzv. Kernel Panic).
    \end{itemize}
    
    \item \textbf{Mikrojádro (např. MINIX, QNX):} V jaderném prostoru běží jen naprosté minimum (plánování CPU, základní paměť a komunikace IPC). Vše ostatní (ovladače, filesystémy) běží odděleně jako běžné uživatelské procesy (v User Space).
    \begin{itemize}
        \item \textit{Výhoda:} Extrémní stabilita a bezpečnost (když spadne ovladač síťové karty, systém nepadne, ovladač se jednoduše restartuje).
        \item \textit{Nevýhoda:} Nižší výkon (velké zdržení a režie kvůli neustálému přepínání kontextu a posílání zpráv mezi procesy).
    \end{itemize}
\end{itemize}

\subsubsection*{Datové struktury jádra}
\label{prg:datastructs}
Jádro je napsáno převážně v jazyce C a ke správě systému využívá datové struktury:
\begin{itemize}
    \item \textbf{\texttt{task\_struct}:} Záznam o procesu (PCB), viz \hyperref[prg:taskstruct]{předchozí kapitola}.
    \item \textbf{\texttt{signal\_struct}:} Datová struktura pro správu signálů.
    \item \textbf{\texttt{inode}:} Datová struktura nesoucí metadata o souboru a ukazatele na datové bloky na disku.
    \item \textbf{\texttt{page}:} Struktura reprezentující jednu fyzickou stránku operační paměti.
\end{itemize}

\subsubsection*{Systém vyrovnávací paměti (Cache / Buffers)}
Operace čtení a zápisu na fyzický disk jsou extrémně pomalé. Jádro proto využívá volnou kapacitu RAM k dočasnému ukládání diskových dat.
\begin{itemize}
    \item \textbf{Page Cache (Mezipaměť stránek):} Ukládá samotný obsah souborů, které byly nedávno čteny nebo zapisovány. Pokud uživatel otevře stejný soubor znovu, načte se z RAM.
    \item \textbf{Buffer Cache:} Slouží k ukládání surových bloků metadat (např. i-uzlů nebo superbloků) z disku. \textit{V moderním Linuxu jsou Page a Buffer cache sjednoceny.}
    \item \textbf{Dirty pages (Špinavé stránky):} Pokud proces zapíše data do souboru, jádro je zapíše pouze do paměti (stránka se označí jako "špinavá"). Fyzický zápis na disk se zpozdí a provádí se hromadně na pozadí (démon \texttt{pdflush}) pro zvýšení výkonu. Vynutit zápis lze příkazem \texttt{sync}.
\end{itemize}

\subsubsection*{Subsystém V/V (Vstup / Výstup)}
\label{prg:vv}
Komunikace s hardwarem je abstrahována tak, aby s ní mohly procesy pracovat jednotně.
Zařízení jsou rozdělená na znaková a bloková viz. \hyperref[prg:typysoub]{typy souborů}.

\begin{itemize}
    \item \textbf{VFS (Virtual File System):} Abstraktní vrstva. Uživatelský program zavolá univerzální systémové volání \texttt{read()} nebo \texttt{write()}. VFS požadavek přijme a předá ho konkrétnímu ovladači (např. pro ext4, NTFS nebo USB flash disk).
    \item \textbf{Přerušení (Interrupts / SysRQ):} Hardwarová zařízení upozorňují procesor na událost (např. stisk klávesy nebo dokončení čtení z disku) pomocí asynchronního přerušení.
    \item \textbf{Struktura termios:} Terminálové rozhraní.
\end{itemize}
Charakteristika zařízení je obsažena v
tabulce deskriptorů zařízení a
zahrnuje:
\begin{itemize}
    \item Identifikátor zařízení
    \item Instrukce zařízení
    \item Stav zařízení
    \item Identifikátor procesu, který se zařízením pracuje
\end{itemize}

\subsubsection*{Správa paměti}
Linux nevyhrazuje procesům pevnou fyzickou paměť, ale využívá systém \textbf{virtuální paměti}. Každý proces má iluzi, že má k dispozici souvislý, obrovský blok vlastní paměti.
\begin{itemize}
    \item \textbf{Stránkování (Paging):} Virtuální paměť i fyzická RAM jsou rozděleny na bloky o stejné velikosti, tzv. \textbf{stránky} (Pages). 
    \item \textbf{MMU a tabulka stránek:} Hardwarová komponenta procesoru \textbf{MMU} (Memory Management Unit) za běhu překládá virtuální adresy na fyzické adresy pomocí Tabulky stránek, kterou spravuje jádro.
    \item \textbf{Výpadek stránky (Page Fault):} Pokud se proces pokusí přistoupit k paměti, která momentálně není nahraná ve fyzické RAM, vznikne výpadek. Jádro proces pozastaví, načte chybějící data z disku do RAM a proces obnoví.
    \item \textbf{Swapping:} Pokud začne docházet fyzická RAM, jádro přesune neaktivní stránky z RAM na pevný disk (do odkládacího oddílu – \textbf{Swap partition}). Jakmile je proces znovu potřebuje, přesunou se zpět. Toto zabraňuje pádu systému, ale výrazně jej zpomaluje.
\end{itemize}

\subsection{Operační systémy reálného času (RTOS)}

\subsubsection*{Definice a základní charakteristika}
\paragraph*{Definice:} RTOS je systém, kde správnost výpočtu nezávisí pouze na logickém výsledku, ale striktně i na čase, kdy je výsledek doručen.
V informatickém pojetí je RTOS
takový OS, který zaručuje, že je
každá správná informace ve správný
čas na správném místě.
\paragraph*{Charakteristika:} Hlavním cílem RTOS není vysoká celková propustnost (jako u běžných OS typu Windows/Linux), ale absolutní predikovatelnost a determinismus. Systém zaručuje, že kritická operace proběhne do přesně stanoveného časového limitu (tzv. \textit{deadline}).

\subsubsection*{Hard a Soft RTOS}
Podle následků zmeškání časového limitu (deadline) se RTOS dělí na:
\begin{itemize}
    \item \textbf{Hard RTOS (Tvrdý):} Zmeškání limitu představuje fatální selhání systému, které může ohrozit životy nebo majetek (např. vystřelení airbagu v autě, řízení kardiostimulátoru, brzdy vlaku).
        Ignoruje data která nepřišla včas.

    \textbf{Point-in-time} - do urřitého času (absolutní)
    \item \textbf{Soft RTOS (Měkký):} Zmeškání limitu nezpůsobí havárii, pouze se sníží kvalita poskytované služby (např. vynechání snímku při streamování videa, VoIP telefonie).
    
    \textbf{Delta-time} - za určitý interval (relativní)
\end{itemize}

\subsubsection*{Hlavní faktory (metriky) RTOS}
\begin{itemize}
    \item Minimální latence při reakci na událost
    \item Minimální latence při přepínání vláken
    \item Někdy nutnost malých rozměrů (embedded systémy)
    \item Minimalizace časových okamžiků, kdy je vypnuto přerušení
    \item Preemptivní plánování založené na prioritách
    \item Dosažení paralelního běhu pokud možno všech zdrojů
    \item Podpora vysoce urgentních úloh
    \item Ochrana před uživateli ztrácí význam (je určen jen pro jednoho uživatele)
    \item Vyžadovaný mechanizmy synchronizace a komunikace mezi úlohami
\end{itemize}

\subsubsection*{Plánování v RTOS (RMS a EDF)}
Běžné OS (např. Linux) využívají plánovače zaměřené na spravedlnost (Fairness). RTOS místo toho využívá algoritmy založené na čase:
\begin{itemize}
    \item \textbf{RMS (Rate Monotonic Scheduling):} Algoritmus se statickou prioritou. Pravidlo je jednoduché: čím kratší má úloha periodu opakování (čím častěji se musí spouštět), tím vyšší fixní prioritu od systému dostane.
    \item \textbf{EDF (Earliest Deadline First):} Algoritmus s dynamickou prioritou. Systém neustále přepočítává priority běžících procesů. Nejvyšší prioritu v daný okamžik dostane ta úloha, jejíž limit pro dokončení (\textit{deadline}) vyprší nejdříve.
\end{itemize}

\subsubsection*{Příklady RTOS}
\begin{itemize}
    \item \textbf{QNX:} Komerční mikrokontrolérový OS, masivně využívaný v automobilovém průmyslu a medicíně.
    \item \textbf{VxWorks:} Komerční RTOS (např. řídil vozítka Curiosity a Perseverance na Marsu).
    \item \textbf{FreeRTOS:} Otevřený systém, dnes de facto standard pro malé mikrokontroléry a IoT zařízení.
    \item \textbf{RT-Linux:} Běžný Linux doplněný o tzv. PREEMPT\_RT patch, který mu dodává vlastnosti RTOS.
\end{itemize}

\section{Řízení procesů UNIXu}
\begin{quote}
    \textbf{vytváření procesu, signály, ukončení procesu, vyvolání jiných procesů, reálné a efektivní UID, změna velikosti procesu, časování procesů, plánování procesů, SysRQ (použití, funkce)}
\end{quote}
Každý proces v systému (kromě procesu \texttt{init/systemd} s PID 1) vzniká jako potomek nějakého již existujícího rodičovského procesu.

\subsection{Vytváření procesu (\texttt{fork})}
\begin{itemize}
    \item \textbf{Systémové volání \texttt{fork()}:} Když běžící proces (rodič) zavolá \texttt{fork()}, jádro OS vytvoří jeho přesnou kopii (potomka).
    \item \textbf{Výsledek:} Potomek dostane vlastní unikátní PID, ale zdědí od rodiče otevřené soubory, proměnné i stav paměti. Od okamžiku volání běží oba procesy nezávisle a souběžně.
    \item \textbf{Návratová hodnota:} Aby program poznal, zda aktuálně běží jako rodič nebo potomek, vrací volání \texttt{fork()} různé hodnoty: 0 se vrátí potomkovi, zatímco rodiči se vrátí PID nového potomka.
    \item \textbf{Copy-on-Write (CoW):} Z důvodu výkonu se fyzická paměť při \texttt{fork()} ve skutečnosti nekopíruje ihned. Oba procesy sdílejí stejnou paměť, a ta se zkopíruje až v okamžiku, kdy se jeden z nich pokusí data změnit.
\end{itemize}

\subsubsection*{Vyvolání jiných procesů (\texttt{exec})}
Volání \texttt{fork()} vytvoří pouze identický klon. Pokud chceme spustit úplně nový program, musíme použít rodinu systémových volání \texttt{exec()} (např. \texttt{execvp}, \texttt{execl}).
\begin{itemize}
    \item \textbf{Princip:} \texttt{exec()} nahradí kompletní paměťový prostor (kód, data, zásobník) aktuálního procesu novým programem nahraným.
    \item \textbf{Zachování PID:} Nový program se spouští v kontextu stávajícího procesu, takže jeho PID zůstává nezměněno.
    \item \textbf{Vzor \texttt{fork} + \texttt{exec}:} Toto je standardní UNIXový způsob spouštění programů. Rodič (např. shell) zavolá \texttt{fork()} (vznikne klon shellu) a tento klon okamžitě zavolá \texttt{exec()} (klon se přepíše na požadovaný program).
\end{itemize}

\subsubsection*{Ukončení procesu (\texttt{exit} a \texttt{wait})}
Proces může skončit normálně (dokončením úkolu) nebo abnormálně (chybou, signálem).
\begin{itemize}
    \item \textbf{Systémové volání \texttt{exit()}:} Proces jím oznamuje jádru, že končí. Předává tzv. návratový kód (0~=~úspěch, nenulová hodnota = chyba). Jádro následně uvolní paměť a zavře soubory procesu.
    \item \textbf{Systémové volání \texttt{wait()}:} Rodičovský proces musí zavolat \texttt{wait()}, aby si od jádra vyzvedl návratový kód svého potomka.
    \begin{itemize}
        \item \textbf{Stav Zombie:} Pokud potomek skončí (\texttt{exit}), ale rodič ještě nezavolal \texttt{wait()}, potomek se stává \textit{zombie}. Nezabírá už RAM ani CPU, ale stále blokuje jedno místo v tabulce procesů.
        \item \textbf{Sirotek (Orphan):} Situace, kdy rodič zemře dříve než potomek. Sirotka převezme proces \texttt{init/systemd}, který po jeho skončení automaticky zavolá \texttt{wait()}.
    \end{itemize}    
\end{itemize}

\subsection{Signály}
Signály jsou formou softwarového přerušení. Slouží k asynchronní komunikaci a řízení procesů (např. upozornění na chybu nebo žádost o ukončení).
Zachicují se pomocí \texttt{trap}.
\begin{itemize}
    \item \textbf{Zpracování signálu:} Když proces obdrží signál, jádro přeruší jeho normální běh. Proces může signál \textit{ignorovat}, \textit{zachytit} (spustit vlastní obslužnou funkci - handler), nebo ponechat \textit{výchozí akci} (typicky okamžité ukončení procesu).
    \item \textbf{Běžné signály:}
    \begin{itemize}
        \item \textbf{\texttt{SIGINT}:} Přerušení z klávesnice (vyvoláno stiskem \texttt{Ctrl+C}). Lze ho zachytit a ignorovat.
        \item \textbf{\texttt{SIGTERM}:} Slušná žádost o ukončení. Proces má čas si uložit data a korektně se ukončit. (Výchozí chování příkazu \texttt{kill}).
        \item \textbf{\texttt{SIGKILL}:} Okamžité a brutální zabití procesu přímo jádrem (\texttt{kill -9}). Proces ho \textit{nemůže} zachytit ani ignorovat.
        \item \textbf{\texttt{SIGSEGV}:} Segmentation fault. Jádro zabíjí proces, protože se pokusil přistoupit do paměti, která mu nepatří.
    \end{itemize}
\end{itemize}

\subsection{Reálné a efektivní UID}
\label{sub:UID}
V UNIXu/Linuxu je bezpečnost a přístup k souborům založen na identitě uživatele (User ID - UID). Každý běžící proces má s sebou spojeno několik takových identifikátorů:
\begin{itemize}
    \item \textbf{Reálné UID (RUID):} Identifikuje skutečného uživatele, který proces spustil. Slouží především pro accounting a k určení, kdo může proces ukončit (např. běžný uživatel může zabít jen procesy se svým vlastním RUID).
    \item \textbf{Efektivní UID (EUID):} Jádro se na něj dívá při ověřování oprávnění (zda má proces právo číst/zapisovat do daného souboru). Většinou je EUID stejné jako RUID.
    \item \textbf{SUID (Set-User-ID) bit:} Speciální oprávnění u spustitelných souborů. Pokud běžný uživatel spustí program se SUID bitem (např. příkaz \texttt{passwd} pro změnu hesla), proces získá EUID vlastníka souboru (což je typicky \texttt{root}). Tím proces získá dočasná administrátorská práva pro zápis do systémových souborů, i když RUID zůstává uživatelské.
    Volá se \texttt{setuid(uid)}.
\end{itemize}
Zpět na \hyperref[sub:identita]{Identita uživatele}.

\subsection{Změna velikosti procesu (Správa paměti)}
Během svého života potřebuje proces často alokovat další paměť pro svá data (tzv. halda - Heap). Zvětšení nebo zmenšení paměťového prostoru procesu se děje na úrovni jádra pomocí systémových volání:
\begin{itemize}
    \item \textbf{Systémová volání \texttt{brk()} a \texttt{sbrk()}:} Mění umístění tzv. \textit{program break} (hranice definující konec datového segmentu/haldy procesu). Posunutím hranice nahoru se halda zvětší, posunutím dolů se paměť vrátí systému.
    \item \textbf{Systémové volání \texttt{mmap()}:} Modernější přístup pro alokaci velkých bloků paměti. Vytváří nová mapování ve virtuálním adresním prostoru procesu (často se používá i pro mapování souborů z disku přímo do RAM).
\end{itemize}
\begin{quote}
        \textbf{Poznámka:} Běžný programátor tato volání nepoužívá přímo, ale využívá knihovní funkce C (např. \texttt{malloc} a \texttt{free}), které na pozadí chytře volají \texttt{brk} nebo \texttt{mmap} podle potřeby.
\end{quote}
\paragraph*{Paměťové segmenty procesu:}
\begin{itemize}
    \item Textový segment - obsahuje kód programu
    \item Datový segment - globální a statické proměnné
    \item Heap (halda) - dynamicky alokovaná paměť
    \item Stack (zásobník) - lokální proměnné
\end{itemize}

\subsection{Časování procesů}
Čas procesoru je rozdělen do epoch. V jedné epoše, každý proces má specifikováno
časové kvantum vypočteno na jejím začátku. V jedné epoše, proces může využívat své časové
kvantum po částech. Epocha končí, když všechny běhu schopné
procesy vyčerpaly svá časová kvanta, kdy jsou
přepočítána časová kvanta všech procesů (ne
jenom běhu schopných) a začne nová epocha. Potomek zdědí základní časové kvantum rodiče.

\subsection{Plánování procesů}
Tři druhy plánování procesů (process scheduling):
\begin{itemize}
    \item \textbf{Krátkodobé:} Výběr, kterému z připravených procesů
bude přidělen procesor, ve všech víceúlohových
systémech.
    \item \textbf{Střednědobé:} Výběr, který
blokovaný nebo připravený proces bude odsunut z
vnitřní paměti na disk, je-li vnitřní paměti nedostatek
(swap out, roll out).
    \item \textbf{Dlouhodobé:} Výběr, která úloha bude spuštěna (aby byl počítač co nejvíce vytížen).
\end{itemize}
Protože procesů je více než jader procesoru, musí jádro (komponenta zvaná Plánovač - \textit{Scheduler}) neustále rozhodovat, komu CPU přidělí.
\begin{itemize}
    \item \textbf{CFS (Completely Fair Scheduler):} Výchozí plánovač moderního Linuxu (je preemptivní). Snaží se o maximální spravedlnost – modeluje ideální multitasking, kde každý proces dostává spravedlivý podíl výkonu procesoru na základě své váhy (priority).
    \item \textbf{Hodnota Nice:} V uživatelském prostoru se priorita neřídí přímo, ale tzv. \textit{slušností} (hodnota \texttt{nice}).
    \begin{itemize}
        \item Rozsah je od \textbf{-20} (nejvyšší priorita, nejméně slušný) do \textbf{+19} (nejnižší priorita, velmi slušný – pustí ostatní procesy před sebe).
        \item Výchozí hodnota je 0. Běžný uživatel může hodnotu \texttt{nice} svému procesu pouze zvýšit. Negativní hodnoty může přidělovat pouze \texttt{root}. K úpravám slouží příkazy \texttt{nice} (při spuštění) a \texttt{renice} (za běhu).
    \end{itemize}
\end{itemize}

\subsection{SysRq}
\textbf{SysRq} je kombinace kláves, která umožňuje uživateli posílat nízkoúrovňové příkazy přímo jádru OS. Funguje to i ve chvíli, kdy systém zamrzne, nereaguje nebo nefunguje grafické rozhraní, protože SysRq je obsluhováno přímo ovladačem klávesnice v jádře.
\begin{itemize}
    \item \textbf{Použití:} Vyvolává se kombinací \texttt{Alt + SysRq + <příkazové písmeno>}. (V Linuxu musí být funkce nejprve povolena, např. zápisem do \texttt{/proc/sys/kernel/sysrq}). (Měl by se zapínat pouze pokud je to nutné hrozí napadení systému.)
    \item \textbf{Funkce a mnemotechnická pomůcka REISUB:} Sekvence pro bezpečné a čisté restartování zamrzlého systému (bez poškození souborového systému):
    \begin{itemize}
        \item \textbf{R} (Raw): Převezme kontrolu nad klávesnicí od zmrzlého grafického serveru (X11/Wayland).
        \item \textbf{E} (tErminate): Pošle všem procesům (kromě \texttt{init}) signál \texttt{SIGTERM} (žádost o bezpečné ukončení).
        \item \textbf{I} (kIll): Pošle signál \texttt{SIGKILL} procesům, které na předchozí krok nereagovaly.
        \item \textbf{S} (Sync): Vynutí okamžitý zápis všech dat z RAM mezipaměti (špinavých stránek) na pevný disk.
        \item \textbf{U} (Unmount): Připojí všechny souborové systémy pouze pro čtení (Read-Only), aby nedošlo k poškození metadat při restartu.
        \item \textbf{B} (reBoot): Provede okamžitý restart systému (obdoba zmáčknutí reset tlačítka na case).
    \end{itemize}
\end{itemize}

\section{Základy běžné uživatelské administrace systému UNIX}
\begin{quote}
    \textbf{vytváření souborů a adresářů, operace se soubory a adresáři, prohledávání systému souborů, identita uživatele, procesu, souboru a jejich změna, přístupová práva a jejich nastavení, přesměrování vstupu a výstupu, propojení příkazů, administrace uživatelů, zálohování, programy pro archivaci a kompresi dat, práce s příkazovými interprety, technologie SMART (význam, použití, vybrané údaje)}
\end{quote}

\subsection{Vytváření a operace se soubory a adresáři}
\begin{itemize}
    \item \textbf{Vytváření:}
    \begin{itemize}
        \item \texttt{touch <soubor>}: Vytvoří nový prázdný soubor nebo zaktualizuje časové razítko existujícího.
        \item \texttt{mkdir <adresář>}: Vytvoří nový adresář (přepínač \texttt{-p} vytvoří celou cestu včetně chybějících nadřazených adresářů).
    \end{itemize}
    \item \textbf{Operace:}
    \begin{itemize}
        \item \texttt{cp <zdroj> <cíl>}: Kopírování souborů. Pro adresáře nutno použít rekurzivně (\texttt{-r}).
        \item \texttt{mv <zdroj> <cíl>}: Přesun souborů nebo jejich přejmenování (v UNIXu jde o stejnou operaci).
        \item \texttt{rm <soubor>}: Smazání souboru. Pro adresáře nutno rekurzivně (\texttt{-r}), pro vynucení bez dotazů (\texttt{-f}).
        \item \texttt{rmdir <adresář>}: Smaže adresář, ale pouze pokud je prázdný.
    \end{itemize}
\end{itemize}

\subsection{Prohledávání systému souborů}
\begin{itemize}
    \item \textbf{\texttt{find}:} Komplexní vyhledávání v reálném čase podle různých kritérií (jméno, velikost, čas modifikace, vlastník). Příklad: \texttt{find /etc -name "*.conf"}.
    \item \textbf{\texttt{locate}:} Velmi rychlé hledání cesty k souboru pomocí předem indexované databáze (nutno aktualizovat pomocí \texttt{updatedb}).
    \item \textbf{\texttt{grep}:} Neprohledává samotný filesystém, ale \textit{obsah} souborů. Hledá řádky odpovídající zadanému regulárnímu výrazu.
    \item \textbf{ls, cd}
\end{itemize}

\subsection{Identita uživatele, procesu, souboru a jejich změna}
\label{sub:identita}
\begin{itemize}
    \item \textbf{Uživatel a proces:} Identifikován pomocí UID (User ID) a GID (Group ID). Proces po spuštění přebírá identitu uživatele (viz dříve \hyperref[sub:UID]{RUID/EUID}). Pro zjištění identity se používají příkazy \texttt{id} nebo \texttt{whoami}.
    \item \textbf{Soubor:} Každý soubor má svého vlastníka (uživatele) a vlastnickou skupinu.
    \item \textbf{Změna identity:}
    \begin{itemize}
        \item \texttt{su <uživatel>}: Přepnutí uživatelského účtu v shellu.
        \item \texttt{sudo <příkaz>}: Jednorázové spuštění příkazu s právy jiného uživatele (typicky \texttt{root}).
        \item \texttt{chown <uživatel>:<skupina> <soubor>}: Změna vlastníka a skupiny souboru.
        \item \texttt{chgrp <skupina> <soubor>}: Změna pouze skupiny.
    \end{itemize}
\end{itemize}

\subsection{Přístupová práva a jejich nastavení}
Přístupová práva se definují pro tři kategorie: \textbf{u} (user/vlastník), \textbf{g} (group/skupina) a \textbf{o} (others/ostatní).
\begin{itemize}
    \item \textbf{Základní práva:} \textbf{r} (read / čtení = 4), \textbf{w} (write / zápis = 2), \textbf{x} (execute / spuštění či vstup do adresáře = 1).
    \item \textbf{\texttt{chmod}:} Nástroj pro změnu práv.
    \begin{itemize}
        \item \textit{Oktalově (číselně):} Sečtením hodnot pro každou kategorii. Příklad: \texttt{chmod 755} (vlastník může vše (4+2+1=7), skupina a ostatní mohou jen číst a spouštět (4+1=5)).
        \item \textit{Symbolicky:} Příklad \texttt{chmod u+x} (přidá vlastníkovi právo spouštět), \texttt{chmod g-w} (odebere skupině právo zápisu).
    \end{itemize}
\end{itemize}

\subsection{Přesměrování vstupu/výstupu a propojení příkazů (Pipes)}
Každý proces v UNIXu má automaticky otevřené 3 standardní proudy dat: \textbf{stdin} (0 - standardní vstup, typicky klávesnice), \textbf{stdout} (1 - standardní výstup, typicky monitor) a \textbf{stderr} (2 - chybový výstup).
\begin{itemize}
    \item \textbf{Přesměrování výstupu:} 
    \begin{itemize}
        \item \texttt{$>$}: Přesměruje stdout do souboru (pokud existuje, přepíše ho).
        \item \texttt{$>>$}: Přesměruje stdout a připojí jej na konec souboru (append).
        \item \texttt{$2>$}: Přesměruje chybový výstup (stderr) do souboru.
    \end{itemize}
    \item \textbf{Přesměrování vstupu (\texttt{<}):} Program místo z klávesnice čte vstup ze specifikovaného souboru.
    \item \textbf{Roura / Pipe (\texttt{|}):} Vezme \texttt{stdout} prvního příkazu a přesměruje ho přímo do \texttt{stdin} druhého příkazu. Umožňuje řetězení jednoduchých nástrojů do složitých operací (např. \texttt{ls -la | grep "txt" | sort}).
\end{itemize}

\subsection{Administrace uživatelů}
Základní systémové soubory (dnes už je neupravujeme ručně, ale přes nástroje):
\begin{itemize}
    \item \texttt{/etc/passwd}: Seznam uživatelů, jejich UID, GID, domovský adresář a výchozí shell.
    \item \texttt{/etc/shadow}: Bezpečně uložená (zahashovaná) hesla uživatelů.
    \item \texttt{/etc/group}: Seznam skupin a jejich členů.
\end{itemize}
\textbf{Administrativní příkazy (vyžadují práva roota):}
\begin{itemize}
    \item \texttt{useradd <jméno>}: Vytvoření nového uživatele.
    \item \texttt{usermod <jméno>}: Modifikace existujícího uživatele (např. změna skupin, shellu).
    \item \texttt{userdel <jméno>}: Smazání uživatele.
    \item \texttt{passwd <jméno>}: Nastavení nebo změna hesla uživatele.
\end{itemize}

\subsection{Zálohování, archivace a komprese dat}
\begin{itemize}
    \item \textbf{Archivace (\texttt{tar}):} Slouží ke spojení mnoha souborů/adresářů do jednoho velkého archivačního souboru (tzv. tarball). Sám o sobě nekomprimuje. (\texttt{tar -cvf} = vytvoření, \texttt{tar -xvf} = rozbalení).
    \item \textbf{Komprese (\texttt{gzip, bzip2, xz}):} Zmenšení velikosti souboru. V UNIXu se typicky kombinuje s tar (přepínače \texttt{-z} pro gzip, \texttt{-j} pro bzip2).
    \item \textbf{Zálohování (\texttt{rsync}):} Velmi mocný nástroj pro synchronizaci a zálohování. Kopíruje pouze změněné části souborů (inkrementální záloha), čímž šetří síťový provoz i čas.
    \item \textbf{\texttt{dd}:} Bitová kopie disků nebo diskových oddílů (low-level záloha).
\end{itemize}

\subsection{Práce s příkazovými interprety (Shell)}
Shell (např. \texttt{bash}, \texttt{zsh}, \texttt{sh}) je textové rozhraní tvořící vrstvu mezi uživatelem a jádrem OS.
\begin{itemize}
    \item \textbf{Skriptování:} Shell obsahuje řídící struktury (cykly, podmínky) a umožňuje automatizaci úloh pomocí shell skriptů.
    \item \textbf{Proměnné prostředí (Environment variables):} Konfigurují běhové prostředí aplikací. Významné jsou např. \texttt{\$PATH} (seznam adresářů, kde systém hledá spustitelné soubory), \texttt{\$USER}, \texttt{\$HOME}. Zpřístupňují se pomocí příkazu \texttt{export}.
    \item \textbf{Konfigurace:} Každý uživatel má vlastní konfigurační soubory (např. \texttt{\~{}/.bashrc} nebo \texttt{\~{}/.profile}), kde lze definovat aliasy (vlastní zkratky příkazů) nebo upravit vzhled promptu.
\end{itemize}

\subsection{Technologie SMART}
\textbf{S.M.A.R.T.} (Self-Monitoring, Analysis and Reporting Technology) je monitorovací systém zabudovaný přímo v pevných discích (HDD/SSD).
\begin{itemize}
    \item \textbf{Význam a použití:} Slouží k predikci selhání disku. Neustále sleduje hardwarové parametry a umožňuje administrátorovi včas vyměnit disk před ztrátou dat. V Linuxu se spravuje pomocí sady \texttt{smartmontools} (konkrétně příkaz \texttt{smartctl}).
    \item \textbf{Klíčové údaje (Atributy):}
    \begin{itemize}
        \item \textit{Reallocated Sector Count:} Počet poškozených sektorů, které disk nahradil záložními. Růst této hodnoty jasně značí odcházející disk.
        \item \textit{Power-On Hours:} Celková doba (v hodinách), po kterou byl disk v provozu.
        \item \textit{Temperature:} Provozní teplota disku (přehřívání drasticky snižuje životnost).
        \item \textit{Spin Retry Count (HDD):} Počet marných pokusů o roztočení ploten (značí mechanický problém).
        \item \textit{Wear Leveling Count (SSD):} Indikátor opotřebení paměťových buněk u SSD.
    \end{itemize}
\end{itemize}

\section{Ochrana dat a informací}
\begin{quote}
    \textbf{kategorie bezpečnostních přístupů (fyzická, informační, kybernetická, síťová, personál bezpečnost atp.), analýza aktiv, rizik a hrozeb, bezpečnostní politika, plán zvládání rizika a protiopatření. Klíčová legislativa a mezinárodně uznávané standardy v oblasti informační a kybernetické bezpečnosti. Malware a ochrana proti škodlivému softwaru. Princip AAA v bezpečnosti, zabezpečení komunikace, zvýšení odolnosti systému a zálohování, možnosti ochrany HW}
\end{quote}

Ochrana dat představuje kompromis mezi mírou zabezpecení a použitelností.

\subsection{Kategorie bezpečnostních přístupů}

Ochrana dat a systémů není pouze záležitostí informačních technologií, ale vyžaduje komplexní přístup. Účinné zabezpečení se skládá z několika vzájemně se doplňujících vrstev (tzv. Defense in Depth). Hlavní kategorie bezpečnostních přístupů zahrnují:

\begin{itemize}
    \item \textbf{Fyzická bezpečnost:} Zajišťuje ochranu samotného hardwaru, datových nosičů, serveroven a budov před neoprávněným fyzickým přístupem, krádeží, poškozením nebo živelnými pohromami. \textit{Prvky: bezpečnostní dveře, kamerové systémy (CCTV), biometrické vstupní systémy, UPS (záložní zdroje), protipožární systémy.}

    \item \textbf{Informační bezpečnost:} Široký koncept chránící data a informace bez ohledu na jejich formu (digitální data, ale i tištěné dokumenty či ústní předání). Její podstatou je udržení tzv. \textbf{CIA triády}:
    \begin{itemize}
        \item \textit{Confidentiality (Důvěrnost):} K datům mají přístup jen oprávněné osoby.
        \item \textit{Integrity (Integrita):} Data nesmí být neoprávněně změněna nebo poškozena.
        \item \textit{Availability (Dostupnost):} Data musí být přístupná, když je uživatel potřebuje.
    \end{itemize}

    \item \textbf{Kybernetická bezpečnost:} Podmnožina informační bezpečnosti, která se úzce zaměřuje na ochranu elektronických dat, počítačových systémů, mobilních zařízení a softwaru před digitálními hrozbami a útoky (např. malware, ransomware, DDoS).

    \item \textbf{Síťová bezpečnost:} Specifická část kybernetické bezpečnosti, jejímž cílem je chránit podnikovou síťovou infrastrukturu a data při jejich přenosu. \textit{Prvky: Firewally, VPN (Virtuální privátní sítě) pro šifrovaný přenos, systémy detekce a prevence průniku (IDS/IPS), oddělování sítí (segmentace a VLAN).}

    \item \textbf{Personální bezpečnost:} Řeší lidský faktor. Zaměřuje se na obranu proti hrozbám zevnitř (insider threats) a útokům založeným na sociálním inženýrství (např. phishing). \textit{Prvky: pravidelná školení uživatelů (Security Awareness), prověřování zaměstnanců, offboarding (odebírání přístupů při odchodu).}

    \item \textbf{Organizační (administrativní) bezpečnost:} Zastřešuje celkové řízení bezpečnosti v organizaci. Zahrnuje tvorbu bezpečnostních politik, směrnic, definování rolí, plány obnovy po havárii (Disaster Recovery Plan) a řízení rizik. Zajišťuje také soulad s legislativou a normami (např. GDPR, ISO 27001, Zákon o kybernetické bezpečnosti).

    \item \textbf{Aplikační bezpečnost:} Týká se zabezpečení dat přímo na úrovni softwaru a databází. Zahrnuje využívání bezpečných metod programování (Secure by Design), penetrační testování, pravidelné aktualizace (patch management) a ochranu proti zranitelnostem (např. SQL injection, XSS). Základním pravidlem zde je uplatňování \textbf{principu nejnižšího privilegia} (uživatel má do systému jen taková práva, která nezbytně potřebuje ke své práci).
\end{itemize}

\subsubsection*{Analýza aktiv, hrozeb a rizik}
Řízení rizik (Risk Management) je základním stavebním kamenem efektivní ochrany dat. Nelze chránit vše na 100~\%, proto organizace musí vědět, kam směřovat své omezené zdroje. Tento proces sestává z následujících kroků:
\begin{itemize}
    \item \textbf{Analýza aktiv:} Identifikace a ohodnocení všeho, co má pro organizaci hodnotu a musí být chráněno. Aktivem nejsou jen data, databáze a servery (informační aktiva), ale také zaměstnanci, fyzické budovy, know-how nebo dobré jméno společnosti (reputace).
    \item \textbf{Analýza hrozeb a zranitelností:} 
    \begin{itemize}
        \item \textit{Hrozba} je událost (nebo útočník), která může aktivum poškodit (např. hacker, ransomware, požár, lidská chyba).
        \item \textit{Zranitelnost} je slabina v systému (např. neaktualizovaný software, slabé heslo, chybějící zámek). Hrozba využívá zranitelnost k útoku na aktivum.
    \end{itemize}
    \item \textbf{Analýza rizik:} Riziko vzniká průnikem hrozby a zranitelnosti. Analýza riziko kvantifikuje nebo kvalifikuje na základě vzorce: \textit{Riziko = Pravděpodobnost hrozby $\times$ Potenciální dopad (škoda na aktivu)}.
\end{itemize}

\subsubsection*{Bezpečnostní politika organizace}
Bezpečnostní politika je vrcholový, strategický dokument schválený nejvyšším vedením organizace. Zastřešuje veškeré snažení o ochranu dat a systémů.
\begin{itemize}
    \item \textbf{Účel a rozsah:} Definuje základní cíle, pravidla a povinnosti pro zajištění bezpečnosti informací. Deklaruje podporu vedení.
    \item \textbf{Abstrakce:} Nezabíhá do detailních technických postupů (neřeší např. konfiguraci konkrétního routeru), ale stanovuje závazný rámec (např. ,,veškerá externí datová komunikace musí být šifrována'').
    \item \textbf{Návazná dokumentace:} Z bezpečnostní politiky vycházejí podrobné směrnice a standardy (např. směrnice pro tvorbu hesel, pravidla pro práci z domova, politika čistého stolu).
    \item \textbf{Závaznost:} Politika je závazná pro všechny zaměstnance a externí spolupracovníky, její porušení bývá sankcionováno na základě pracovněprávních předpisů.
\end{itemize}

\subsubsection*{Plán zvládání rizika a protiopatření}
Na základě výsledků analýzy rizik rozhoduje vedení organizace o způsobu, jakým s jednotlivými riziky naloží. Existují čtyři základní strategie:
\begin{itemize}
    \item \textbf{Snížení (Mitigace) rizika:} Nasazení protiopatření, která sníží pravděpodobnost výskytu hrozby nebo minimalizují její dopad.
    \item \textbf{Přenos (Transfer) rizika:} Přenesení finančních či provozních následků na třetí stranu (např. sjednání kybernetického pojištění nebo outsourcing rizikové služby do cloudu).
    \item \textbf{Vyhnutí se riziku:} Úplné ukončení činnosti nebo procesu, který riziko generuje (např. zákaz používání USB flash disků ve firmě, vypnutí zastaralého systému).
    \item \textbf{Akceptace rizika:} Náklady na protiopatření by byly vyšší než hodnota samotného aktiva nebo potenciální škoda. Vedení riziko formálně přijme a je připraveno nést případné následky.
\end{itemize}
\textbf{Protiopatření:} Jsou konkrétní mechanismy (bezpečnostní kontroly) k mitigaci rizik. Dělí se dle účelu na \textit{preventivní} (zabraňují incidentu, např. firewall), \textit{detekční} (upozorní na probíhající incident, např. alarm) a \textit{korekční/reaktivní} (obnova dat ze zálohy po havárii).

\subsection{Legislativa a standardy}

Oblast ochrany dat a kybernetické bezpečnosti je přísně regulována. Zatímco legislativa definuje \textbf{co} musí být chráněno (závazné právní mantinely), mezinárodní standardy poskytují osvědčené postupy (Best Practices), \textbf{jak} toho technicky a procesně dosáhnout.

\subsubsection*{Klíčová legislativa Evropské unie}
\begin{itemize}
    \item \textbf{GDPR (Obecné nařízení o ochraně osobních údajů)} 
    \item \textbf{Směrnice NIS 2 (Network and Information Security)}
    \item \textbf{DORA (Digital Operational Resilience Act)}
\end{itemize}
\subsubsection*{Klíčová legislativa České republiky}
\begin{itemize}
    \item \textbf{Zákon o kybernetické bezpečnosti (ZKB):} Implementuje evropskou směrnici NIS do českého práva. Vymezuje práva a povinnosti pro orgány státní správy, provozovatele kritické informační infrastruktury (KII) a poskytovatele základních služeb. Kontrolním a metodickým orgánem je \textbf{NÚKIB} (Národní úřad pro kybernetickou a informační bezpečnost).
    \item \textbf{Zákon o zpracování osobních údajů:} Adaptační zákon, který doplňuje a upřesňuje pravidla GDPR pro prostředí ČR. Dozorovým orgánem pro tuto oblast je \textbf{ÚOOÚ} (Úřad pro ochranu osobních údajů).
\end{itemize}

\subsubsection*{Mezinárodně uznávané standardy}
Kromě legislativy se organizace řídí normami, které umožňují systematické řízení bezpečnosti. Často slouží jako podklad pro získání certifikace, kterou firmy vyžadují v rámci obchodních smluv.
Organizace vytvářenící normi: ISO, IEC, ITU (v ČR: ČSNI, pro EU: ENISA)
\begin{itemize}
    \item \textbf{Rodina norem ISO/IEC 27000:} Nejznámější globální standardy pro informační bezpečnost.
    \begin{itemize}
        \item \textbf{ISO/IEC 27001:} Specifikuje požadavky na zavedení a udržování \textbf{ISMS} (Information Security Management System - Systém řízení bezpečnosti informací). Je založen na neustálém zlepšování pomocí Demingova cyklu \textbf{PDCA} (Plan - Do - Check - Act). Organizace se dle této normy mohou nechat certifikovat.
        \item \textbf{ISO/IEC 27002:} Obsahuje podrobný katalog konkrétních bezpečnostních opatření (kontrol), které pomáhají naplnit požadavky normy 27001.
    \end{itemize}
    \item \textbf{NIST Cybersecurity Framework (CSF):} Velmi populární rámec vyvinutý americkým Národním institutem standardů a technologie. Dělí obranu do pěti logických fází, které se dobře komunikují vedení firmy: \textbf{Identify} (Identifikuj aktiva a rizika) $\rightarrow$ \textbf{Protect} (Chraň je) $\rightarrow$ \textbf{Detect} (Detekuj útok) $\rightarrow$ \textbf{Respond} (Reaguj na incident) $\rightarrow$ \textbf{Recover} (Obnov provoz).
\end{itemize}

\subsection{Malware}

\textbf{Malware} (z anglického \textit{malicious software} – škodlivý software) je zastřešující pojem pro jakýkoliv program nebo kód vytvořený s cílem poškodit počítačový systém, odcizit data, získat neoprávněný přístup nebo vydírat oběť. 

\subsubsection*{Základní typy malwaru}
Škodlivý kód se klasifikuje především podle způsobu šíření a svého chování:
\begin{itemize}
    \item \textbf{Počítačový virus:} Připojuje se k legitimním spustitelným souborům. K aktivaci a šíření vyžaduje interakci uživatele (např. spuštění infikovaného souboru).
    \item \textbf{Červ (Worm):} Na rozdíl od viru je samostatný a nepotřebuje hostitelský soubor ani interakci uživatele. K masivnímu šíření sítí využívá bezpečnostní zranitelnosti OS a aplikací.
    \item \textbf{Trojský kůň (Trojan):} Maskuje se jako užitečný nebo legitimní program. Sám se nereplikuje, ale po instalaci uživatelem často otevírá zadní vrátka do systému (Backdoor) pro stažení dalšího malwaru.
    \item \textbf{Ransomware:} Zašifruje data na disku (nebo uzamkne zařízení) a za dešifrovací klíč požaduje výkupné. Moderní varianty data před zašifrováním navíc ukradnou a hrozí jejich zveřejněním.
    \item \textbf{Spyware a Keyloggery:} Skrytě sledují činnost uživatele, kradou přihlašovací údaje, historii prohlížení a zaznamenávají stisky kláves.
    \item \textbf{Rootkit:} Sofistikovaný malware operující hluboko v systému (často na úrovni jádra OS). Jeho cílem je utajit svou přítomnost a přítomnost dalšího škodlivého kódu před antivirem.
    \item \textbf{Botnet:} Síť infikovaných zařízení (zvaných zombie), která útočník ovládá na dálku. Slouží k rozesílání spamu nebo k hromadným \textbf{DDoS} útokům.
\end{itemize}

\subsubsection*{Ochrana a protiopatření}
Obrana proti malwaru musí být vícevrstvá, nelze spoléhat pouze na jeden nástroj:
\begin{itemize}
    \item \textbf{Antivirové systémy a EDR:} Běžný antivirus funguje na bázi detekce \textit{virových signatur} (známých vzorů v kódu). Moderní firemní ochrana využívá \textbf{EDR} (Endpoint Detection and Response), která analyzuje nestandardní \textit{chování} procesů v reálném čase pomocí AI a dokáže zastavit i dosud neznámý malware (Zero-day hrozby).
    \item \textbf{Patch Management (Aktualizace):} Pravidelné a včasné instalování bezpečnostních záplat operačního systému a aplikací, což eliminuje zranitelnosti, které využívají červi.
    \item \textbf{Filtrování sítě a e-mailů:} Většina malwaru (zejména ransomware a trojské koně) přichází formou \textbf{phishingu} v e-mailu. Nutností je blokování nebezpečných příloh a zakázání spouštění maker v dokumentech MS Office.
    \item \textbf{Zálohování:} Jediná 100\% spolehlivá obrana proti ransomwaru. Ideálně dle pravidla \textbf{3-2-1} (3 kopie dat, na 2 různých médiích, z toho 1 geograficky oddělena/offline). Zálohy navíc musí být chráněny proti smazání (tzv. immutable backups).
    \item \textbf{Vzdělávání uživatelů:} Školení zaměstnanců např. proti phishingu.
\end{itemize}

\subsection{Princip AAA v bezpečnosti}
Proces řízení přístupu k datům a systémům se opírá o bezpečnostní model \textbf{AAA}, který se skládá ze tří na sebe navazujících kroků (předchází jim \textit{Identifikace} – uživatel řekne, kdo je, např. zadá login):
\begin{itemize}
    \item \textbf{Autentizace (Authentication – Ověření totožnosti):} Uživatel dokazuje, že je skutečně tím, za koho se vydává. Provádí se pomocí něčeho, co uživatel \textit{zná} (heslo, PIN), co \textit{má} (čipová karta, token, mobil) nebo čím \textit{je} (biometrie – otisk, Face ID). Dnes je standardem \textbf{MFA} (Vícefaktorová autentizace).
    \item \textbf{Autorizace (Authorization – Oprávnění):} Systém na základě ověřené identity přidělí uživateli přístupová práva. Využívá se zde \textbf{RBAC} (Role-Based Access Control) – práva se nepřidělují jednotlivcům, ale pracovním rolím. Platí princip \textbf{nejnižšího privilegia} (přístup jen k tomu, co je nezbytné pro práci).
    \item \textbf{Auditování (Accounting – Záznam):} Monitorování a logování toho, co uživatel v systému dělal (kdy se přihlásil, jaké soubory smazal). Slouží k vyšetřování incidentů a zajištění principu \textit{nepopiratelnosti} (Non-repudiation).
\end{itemize}

\subsection{Zabezpečení komunikace}
Cílem je ochrana dat při jejich přenosu (Data in Transit) před odposlechem a modifikací. Zajišťuje se pomocí kryptografie a síťových protokolů:
\begin{itemize}
    \item \textbf{Šifrování sítě:} Pro bezpečný přístup do firemní sítě přes nezabezpečený internet se využívá \textbf{VPN} (Virtuální privátní síť), která vytváří šifrovaný tunel (protokoly IPsec nebo OpenVPN).
    \item \textbf{Zabezpečení aplikační vrstvy:} Přenos webových stránek a API komunikace musí probíhat přes \textbf{HTTPS/TLS}. Pro e-maily či messengery se využívá \textbf{E2EE} (End-to-End šifrování), kdy data dokáže dešifrovat pouze odesílatel a příjemce (ani provozovatel služby k nim nemá přístup).
    \item \textbf{Digitální podpisy a haše:} Pro ověření integrity zpráv (zda nebyly po cestě změněny) a ověření identity odesílatele (asymetrická kryptografie, infrastruktura veřejných klíčů -- PKI).
\end{itemize}

\subsection{Zvýšení odolnosti systému a zálohování}
Odolnost znamená schopnost systému pokračovat v provozu i při selhání hardwaru nebo útoku.
\begin{itemize}
    \item \textbf{Vysoká dostupnost (High Availability) a Redundance:} Kritické prvky (servery, disky, switche, zdroje napájení) jsou zdvojeny. Při výpadku primárního prvku automaticky přebírá zátěž záložní (Failover). Využívají se disková pole \textbf{RAID} a clustering serverů.
    \item \textbf{Zálohování (Backup):} Vytváření kopií dat pro případ ztráty. Definuje se dvěma klíčovými metrikami:
    \begin{itemize}
        \item \textbf{RPO (Recovery Point Objective):} O kolik dat si může firma dovolit přijít (např. zálohujeme každou hodinu = RPO je 1 hodina).
        \item \textbf{RTO (Recovery Time Objective):} Za jak dlouho po havárii musí být systém opět plně funkční.
    \end{itemize}
    \item \textbf{Typy záloh:} \textit{Plná} (všechna data, pomalá), \textit{Inkrementální} (pouze změny od poslední jakékoliv zálohy, rychlá tvorba, pomalá obnova), \textit{Diferenciální} (změny od poslední plné zálohy).
\end{itemize}

\subsection{Možnosti ochrany hardwaru}
Kromě klasické fyzické bezpečnosti (zámky, mříže, kamery, UPS, generátory, biometrický přístup do serveroven) se HW chrání i na úrovni samotného zařízení:
\begin{itemize}
    \item \textbf{Šifrování celého disku (FDE - Full Disk Encryption):} Nástroje jako BitLocker (Windows) nebo FileVault (macOS) šifrují pevný disk. Pokud je notebook nebo disk fyzicky ukraden, data jsou bez dešifrovacího klíče nečitelná.
    \item \textbf{Čip TPM (Trusted Platform Module):} Specializovaný kryptografický čip na základní desce zařízení. Bezpečně uchovává šifrovací klíče a zajišťuje proces tzv. \textit{Secure Boot} (ověřuje, že se při startu počítače nespouští škodlivý kód, např. rootkit).
    \item \textbf{Ochrana portů a BIOS/UEFI:} Ochrana vstupu do BIOSu heslem, zamezení bootování z externích médií a softwarové či fyzické uzamčení USB portů proti vložení neoprávněných flash disků (ochrana před únikem dat – DLP, i před zavlečením malwaru).
\end{itemize}

\section{Kryptografie}
\begin{quote}
    \textbf{základní principy a pojmy v oblasti kryptografie, klasifikace šifer (substituční, monoalfabetické, homofonní atp.), kryptografický systém, klíč, časová a paměťová náročnost, kryptoanalýza. Hashovací funkce, steganografie, symetrické a asymetrické šifrování, infrastruktura veřejných klíčů, certifikáty, certifikační autority}
\end{quote}

\subsection{Základní principy a pojmy v oblasti kryptografie}

Kryptografie je vědní disciplína (součást kryptologie), která se zabývá tvorbou a zkoumáním metod pro bezpečný přenos a ukládání dat za přítomnosti protivníka (útočníka). Její podstatou je transformace srozumitelných dat do nesrozumitelné podoby a naopak s využitím matematických aparátů.

\subsubsection*{Základní pojmy}
\begin{itemize}
    \item \textbf{Otevřený text (Plaintext):} Původní, srozumitelná zpráva nebo data určená k zašifrování.
    \item \textbf{Šifrový text (Ciphertext):} Zašifrovaná, pro neoprávněnou osobu nesrozumitelná zpráva.
    \item \textbf{Šifrování (Encryption):} Proces transformace otevřeného textu na šifrový text pomocí šifrovacího algoritmu a klíče.
    \item \textbf{Dešifrování (Decryption):} Opačný proces, kdy oprávněný příjemce transformuje šifrový text zpět na otevřený text.
    \item \textbf{Klíč (Key):} Tajný parametr (sekvence bitů), který řídí proces šifrování a dešifrování. Zajišťuje, že i při znalosti šifrovacího algoritmu nelze zprávu bez správného klíče přečíst.
    \item \textbf{Šifra / Šifrovací algoritmus:} Matematická funkce (nebo sada instrukcí), která definuje, jak přesně probíhá transformace dat.
\end{itemize}

\subsubsection*{Cíle kryptografie}
Kryptografické mechanismy neslouží pouze k samotnému utajení obsahu zprávy, ale plní čtyři základní funkce moderní informační bezpečnosti:
\begin{itemize}
    \item \textbf{Důvěrnost (Confidentiality):} Zajištění, že data si může přečíst pouze ten uživatel, který vlastní správný dešifrovací klíč.
    \item \textbf{Integrita (Integrity):} Schopnost spolehlivě detekovat, zda data nebyla během přenosu nebo uložení neoprávněně (či náhodně) změněna. K tomuto účelu slouží tzv. hašovací funkce.
    \item \textbf{Autentizace (Authentication):} Prokázání identity komunikující strany, tedy bezpečné ověření, že data skutečně pocházejí od deklarovaného odesílatele.
    \item \textbf{Nepopiratelnost (Non-repudiation):} Odesílatel zprávy nemůže v budoucnu lživě popřít, že zprávu odeslal nebo vytvořil. Typicky se realizuje pomocí asymetrické kryptografie a digitálních podpisů.
\end{itemize}

\subsubsection*{Kerckhoffsův princip}
Zásadní teoretické pravidlo moderní kryptografie (formulované v 19. století), které říká, že \textbf{bezpečnost šifrovacího systému nesmí záviset na utajení samotného algoritmu, ale pouze na utajení klíče}. 

\subsection{Klasifikace šifer (Základní operace a klasické šifry)}

Z hlediska základního principu fungování dělíme šifry do dvou hlavních kategorií: \textbf{transpoziční} a \textbf{substituční}.

\subsubsection*{1. Transpoziční (permutační) šifry}
Znaky otevřeného textu se nenahrazují jinými, ale pouze se mění jejich pořadí (přeskupují se) podle určitého pravidla. Z textu se stává přesmyčka (anagram).
Cílem je difúze informacního obsahu zprávy po celé její délce.
Kryptoanalýza je namáhavejší než u substitucních šifer.
Nevýhodou použití jsou pametové a casové nároky pri šifrování.
\begin{itemize}
    \item \textit{Příklad:} Skytalé (starověké Řecko - navíjení pergamenu na válec určitého průměru), Rail Fence (psaní textu do "plotu").
\end{itemize}

\subsubsection*{2. Substituční šifry}
Každý znak (nebo skupina znaků) otevřeného textu je nahrazen jiným znakem, symbolem nebo číslem podle stanoveného klíče. Původní pořadí znaků zůstává zachováno. Dále se dělí na několik typů:

\begin{itemize}
    \item \textbf{Monoalfabetické šifry:} K šifrování se používá pouze jedna šifrovací abeceda. Jeden konkrétní znak otevřeného textu je vždy nahrazen stejným znakem šifrového textu (např. každé "A" se změní na "D").
    \begin{itemize}
        \item \textit{Příklad:} Caesarova šifra (posun abecedy o pevný počet míst).
        \item \textit{Zranitelnost:} Velmi snadno prolomitelné pomocí \textbf{frekvenční analýzy} (porovnání četnosti znaků v šifrovém textu s průměrnou četností znaků v daném přirozeném jazyce).
    \end{itemize}

    \item \textbf{Homofonní šifry:} Byly vytvořeny jako obrana proti frekvenční analýze. Jeden znak otevřeného textu může být nahrazen několika různými znaky (nebo čísly) šifrového textu. Znakům, které se v jazyce vyskytují často (např. samohláska "E"), je přiřazeno více možných šifrových variant, čímž se vyrovná rozložení četnosti v šifrovém textu.

    \item \textbf{Polygramové šifry:} Nenahrazují znaky po jednom (po písmenech), ale šifrují skupiny znaků najednou (např. dvojice písmen - bigramy, nebo trojice - trigramy).
    \begin{itemize}
        \item \textit{Příklad:} Playfairova šifra (šifruje dvojice písmen pomocí matice 5x5).
    \end{itemize}

    \item \textbf{Polyalfabetické šifry:} Používají více šifrovacích abeced současně. To, jakou abecedou se daný znak zašifruje, závisí na jeho pozici v textu a na použitém klíči (hesle). Jedno písmeno "A" tak může být v šifrovém textu reprezentováno pokaždé jiným znakem.
    \begin{itemize}
        \item \textit{Příklad:} Vigenèrova šifra. Odolnost proti běžné frekvenční analýze je zde mnohem vyšší.
    \end{itemize}
\end{itemize}

\subsection{Kryptografický systém (Kryptosystém)}
Kryptosystém je komplexní soubor prvků, algoritmů a postupů, které společně umožňují bezpečné šifrování a dešifrování dat. Podle způsobu šifrování a přenášení klíče se rozděluje na varianty:
\begin{itemize}
    \item Symetrické
    \item Asymetrické
    \item Hybridní
\end{itemize}

\subsection{Klíč}
Klíč je tajný parametr, který vstupuje do šifrovacího či dešifrovacího algoritmu a mění jeho chování. V souladu s Kerckhoffsovým principem je bezpečnost celého systému závislá výhradně na utajení tohoto klíče.
\begin{itemize}
    \item \textbf{Délka klíče:} Udává se v bitech a určuje velikost prostoru klíčů (např. klíč o délce $n$ bitů má $2^n$ možných kombinací). Čím delší klíč je, tím obtížnější je jeho prolomení.
\end{itemize}

\subsection{Kryptoanalýza}
Kryptoanalýza je vědní disciplína zabývající se hledáním slabin v kryptografických systémech a jejich prolamováním bez předchozí znalosti klíče. Podle informací, které má útočník k dispozici, rozlišujeme základní modely útoků:
\begin{itemize}
    \item \textbf{Útok pouze se znalostí šifrového textu (Ciphertext-only):} Útočník má k dispozici pouze zašifrované zprávy a snaží se z nich odvodit otevřený text nebo klíč.
    \item \textbf{Útok se znalostí otevřeného textu (Known-plaintext):} Útočník zná (nebo uhodne) textovou podobu části zprávy a zároveň má k dispozici její šifrovaný ekvivalent.
    \item \textbf{Útok s volbou otevřeného textu (Chosen-plaintext):} Útočník má možnost nechat zašifrovat libovolný vlastní text a analyzovat vzniklý šifrový text.
    \item \textbf{Útok hrubou silou (Brute-force):} Systematické zkoušení všech možných kombinací klíčů z prostoru $K$, dokud se neobjeví smysluplný otevřený text.
    \item \textbf{Další:} Slovníkoví útok, Frekvenční analýza
\end{itemize}

\subsection{Výpočetní náročnost}
Bezpečnost moderních šifer se opírá o tzv. výpočetní bezpečnost — útok je teoreticky možný, ale prakticky neproveditelný z důvodu limitovaných zdrojů (času a paměti). Hodnotí se proto asymetricky:
\begin{itemize}
    \item \textbf{Časová náročnost:} Vyjadřuje množství operací (procesorového času) potřebných k provedení algoritmu. Pro legitimní uživatele musí být šifrování a dešifrování rychlé (polynomiální čas $\mathcal{O}(n^k)$). Pro útočníka bez klíče musí být kryptoanalýza časově nezvládnutelná (exponenciální čas $\mathcal{O}(2^n)$), trvající například miliony let i na superpočítačích.
    \item \textbf{Paměťová náročnost:} Vyjadřuje množství operační paměti (RAM) nebo úložného prostoru, které algoritmus během svého běhu vyžaduje. Většina běžných šifer je navržena s ohledem na minimální paměťovou náročnost, aby fungovaly na vestavěných zařízeních či čipových kartách. Naopak u specifických funkcí (např. hashování hesel jako Argon2) se paměťová náročnost záměrně zvyšuje, což slouží jako obrana proti hardwarové akceleraci útoků pomocí GPU a ASIC čipů.
\end{itemize}

\subsection{Hashovací funkce}
Matematické algoritmy, které přijímají vstupní data libovolné délky a převádějí je na výstup pevné, předem dané délky (tzv. hash, otisk zprávy nebo fingerprint).
\begin{itemize}
    \item \textbf{Základní vlastnosti:}
    \begin{itemize}
        \item \textbf{Jednosměrnost:} Ze získaného hashe nelze matematicky zkonstruovat původní zprávu.
        \item \textbf{Bezkoliznost:} Je výpočetně neproveditelné najít dvě různé zprávy, které by měly stejný hash.
        \item \textbf{Citlivost na vstup:} I nepatrná změna na vstupu (např. změna jednoho znaku) vede ke zcela odlišnému výslednému hashi.
    \end{itemize}
    \item \textbf{Využití:} Kontrola integrity dat (zda nebyl soubor poškozen či pozměněn), bezpečné ukládání hesel (neukládá se heslo, ale pouze jeho hash) a tvorba digitálních podpisů. Příklady: rodina SHA-2, SHA-3.
\end{itemize}

\subsection{Steganografie}
Na rozdíl od kryptografie, která zprávu činí nesrozumitelnou (ale je zjevné, že jde o šifru), steganografie skrývá samotnou existenci zprávy.
\begin{itemize}
    \item \textbf{Princip:} Tajná data jsou vložena do jiného, na první pohled neškodného nosiče (tzv. cover medium), aby nevzbudila pozornost útočníka.
    \item \textbf{Využití:} Úprava nejméně významných bitů (LSB) v obrazových (JPEG, PNG) nebo zvukových souborech. Často se kombinuje s kryptografií — zpráva se nejprve bezpečně zašifruje a následně se vzniklý šifrový text steganograficky skryje.
\end{itemize}

\subsection{Symetrické a asymetrické šifrování}
Liší přístupem ke kryptografickým klíčům:
\begin{itemize}
    \item \textbf{Symetrické šifrování:} Používá stejný klíč pro šifrování i dešifrování.
    \begin{itemize}
        \item \textit{Výhody:} Velmi rychlé a výpočetně nenáročné, ideální pro šifrování velkých objemů dat.
        \item \textit{Nevýhody:} Zásadní problém s distribucí klíče — obě komunikující strany si musí před zahájením komunikace bezpečně předat tajný klíč tak, aby jej nikdo neodposlechl.
        \item \textit{Příklady:} 
        \begin{itemize}
            \item blokové -- AES (Advanced Encryption Standard), DES, ChaCha20, RC5, Blowfish
            \item Proudové -- RC4
        \end{itemize}
    \end{itemize}
    \item \textbf{Asymetrické šifrování:} Používá matematicky provázaný pár klíčů — veřejný klíč (Public Key) a soukromý klíč (Private Key). Co jedním klíčem zašifrujeme, to lze dešifrovat pouze druhým klíčem z páru.
    \begin{itemize}
        \item \textit{Princip:} Veřejný klíč odesílatel volně zveřejní. Kdokoliv jím může zašifrovat zprávu, ale dešifrovat ji může pouze majitel patřičného soukromého klíče, který drží v tajnosti.
        \item \textit{Výhody:} Řeší problém bezpečné distribuce klíčů a umožňuje realizaci digitálních podpisů (zajišťuje nepopiratelnost a autentizaci).
        \item \textit{Nevýhody:} Výpočetně velmi náročné a pomalé (běžně se proto používá k bezpečné výměně symetrického klíče, kterým se pak šifrují samotná data — tzv. hybridní šifrování).
        \item \textit{Příklady:} Diffie-Hellman, El Gamal, RSA, kryptografie eliptických křivek (ECC), DSA pro elektronické podpisy.
    \end{itemize}
\end{itemize}

\subsection{Infrastruktura veřejných klíčů (PKI) a certifikáty}
Aby asymetrická kryptografie fungovala v praxi (např. na internetu), je nutné zajistit, že daný veřejný klíč skutečně patří tomu, kdo to o sobě tvrdí (obrana proti útoku Man-in-the-Middle). K tomu slouží systém PKI.
\begin{itemize}
    \item \textbf{PKI (Public Key Infrastructure):} Ucelený systém hardwaru, softwaru, lidí, politik a postupů potřebných pro vytváření, správu, distribuci a revokaci digitálních certifikátů.
    \item \textbf{Digitální certifikát:} Elektronický dokument (typicky ve formátu X.509), který pevně svazuje veřejný klíč s identitou jeho majitele (např. s doménou webového serveru nebo jménem osoby).
    
    Obsahuje:
    \begin{itemize}
        \item Veřejný klíč
        \item Informace o držiteli
        \item Platnost certifikátu
        \item Název a podpis certifikační autority
    \end{itemize}
    \item \textbf{Certifikační autorita (CA):} Důvěryhodná třetí strana, která funguje jako digitální notář. CA ověří identitu žadatele a pokud je vše v pořádku, žadatelův certifikát (obsahující jeho veřejný klíč) elektronicky podepíše svým vlastním soukromým klíčem. Pokud uživatelský počítač důvěřuje této CA, bude automaticky důvěřovat i všem certifikátům, které vydala. (např. Lets encrypt, DigiCert, Globalsign, Comodo)
\end{itemize}